\selectlanguage{english}

In this chapter, we develop the concept of black hole (BH). First, black holes need to be rigorously defined (\secref{sec:bh-def}): in order to do that, it is necessary to study the global structure of spacetime, and in particular its asymptotics, with a techinque developed by Penrose \cite{penrose-1, penrose-2} know as \emph{conformal completion}. Secondly, we perform a geometric analysis of event horizons (\secref{sec:event-horizons}), in order to study properties of black holes which are independent of the gravitational theory used. Then, we study black holes at equilibrium (sec:bh-stationary), covering the three main stationary black holes solutions: Schwarschild (\secref{sec:bh-schwarzschild}), Reissner--Nordström (\secref{sec:bh-rn}) and Kerr black holes (\secref{sec:bh-kerr}). Finally, we study the dynamics of black holes out of equilibrium (\secref{sec:bh-dynamics}).

\section{Definition of black hole}
\label{sec:bh-def}

\subsection{Naive definition}

It is possible to give a naive and intuitive definition of BH: given a spacetime $ (\man , g) $, a \bctxt{black hole} is ``a localized region of spacetime from which neither massive nor massless particles can escape". This definition has two essential features: \emph{localization} and \emph{inescapability}. The latter implies the existence of a boundary which no particle emitted in the black hole can cross: this boundary is known as the \bctxt{event horizon} (or simply horizon), and it is a \emph{one-way membrane}, in the sense that it can be crossed from the black hole's ``exterior" towards its ``interior", but not in the reverse way. This means that the horizon must be a hypersurface which divides $ \man $ int two regions: the interior (i.e. the BH itself) and the exterior of the black hole.

\subsubsection{Black hole boundaries}
\label{sssec:bh-boundaries}

Recall that, on a Lorentzian $ d $-manifold $ (\man , g) $, a hypersurface $ \Sigma $ can be locally classified based on its \emph{causal type}, which is determined by the signature of the restricted metric $ g \vert_\Sigma $. The hypersurface $ \Sigma $ is said to be:
\begin{itemize}
  \item \bctxt{spacelike} if $ g \vert_\Sigma $ is positive definite, i.e. if its signature is $ (0 , d-1) $, which means that $ (\Sigma , g \vert_\Sigma) $ is a Riemannian manifold;
  \item \bctxt{timelike} if $ g \vert_\Sigma $ is a Lorentzian metric, i.e. if its signature is $ (1 , d-2) $, which means that $ (\Sigma , g \vert_\Sigma) $ is a Lorentzian manifold;
  \item \bctxt{null} if $ g \vert_\Sigma $ is degenerate, i.e. if its local matrix representation has a vanishing eigenvalue.
\end{itemize}
All these definitions can be applied locally, i.e. to each point of $ p \in \Sigma $, and it is of course possible that $ \Sigma $ has a different causal nature at different points. The causal type of $ \Sigma $ at $ p \in \Sigma $ can also be deduced from any normal vector $ n_p \in N_p \Sigma $, since the causal type of $ \Sigma $ at $ p $ is the opposite of the causal type of $ n_p $, i.e.:
\begin{itemize}
  \item $ \Sigma $ spacelike at $ p $ $ \iff $ $ n_p $ timelike;
  \item $ \Sigma $ timelike at $ p $ $ \iff $ $ n_p $ spacelike;
  \item $ \Sigma $ null at $ p $ $ \iff $ $ n_p $ null;
\end{itemize}
This can be proved by considering an orthogonal basis adapted to $ \Sigma $. Note that null hypersurfaces have the distinctive feature that their normals are also tangent to them, as $ \braket{n , n}_g = 0 $ by definition of null vector.

\begin{figure}
  \centering
  \includegraphics[width = 0.30 \textwidth]{space-membrane.png}
  \hspace{10pt}
  \includegraphics[width = 0.30 \textwidth]{time-membrane.png}
  \hspace{10pt}
  \includegraphics[width = 0.30 \textwidth]{null-membrane.png}
  \caption{The three causal types of hypersurfaces: spacelike (left), timelike (middle) and null (right). Figure from \cite{gourgoulhon}.}
  \label{fig:causal-hyp}
\end{figure}

Clearly, a timelike hypersurface is a two-way membrane: it locally divides spacetime in two regions, and a future-directed timelike or null worldline can cross it in both ways, as shown in \figref{fig:causal-hyp}. On the contrary, a spacelike hypersurface is a one-way membrane: a future-directed timelike or null worldlike, which is constrained to move inside light cones, can cross it only in the future-directed way. A null hypersurface is also a one-way membrane, since, at most, a null worldline which is not going from the past region to the future region must remain on the hypersurface, as depicted in \figref{fig:causal-hyp}. Therefore, a black hole boundary can in principle be either a null or a spacelike hypersurface: it turns out that, wherever it is smooth, the event horizon is a null hypersurface (\pref{prop:achronal-null}).

\subsubsection{Localization}

Inescapability is only one of the two naive defining features of black holes, and it alone is not enough to characterize them. Indeed, any region bounded by a null hypersurface cannot be qualified as a ``black hole": for instance, the interior of a future light cone in Minkowski spacetime is such a region. What distinguishes it from a black hole is that any particle moving inside it, although trapped by the light-cone boundary, can travel infinitely far from any ``center" in that region.

The necessity of the localization condition is then evident, since it forces the black hole region to have a final spatial extent, so that particles inside it cannot travel to infinitely remote regions. To formalize this concept in a coordinate-independent way, one introduces a so-called \emph{conformal completion} of the spacetime $ (\man , g) $, which brings its infinities to a finite distance in another manifold.

\subsection{Conformal completions}

The machinery of conformal completions is based on conformal transformations between spacetimes, which turn out to preserve their causal structure, thus allowing to study their ``behaviour at infinity" and, eventually, to geive a formal definition of black hole.

\subsubsection{Conformal transformations}

\begin{definition}{Conformal embedding}{}
  Let $ (\man , g) $ and $ (\tilde{\man} , \tilde{g}) $ be two metric manifolds. A \bcdef{conformal embedding} of $ \man $ into $ \tilde{\man} $ is a smooth embedding $ \Phi : \man \ra \tilde{\man} $ equipped with $ \Omega \in \cfk{\infty}(\man , \R^+) $ such that:
  \begin{equation}
    \tilde{g} \vert_{\Phi(\man)} = \Omega^2 g
    \label{eq:conformal-relation}
  \end{equation}
\end{definition}

Usually, $ \Phi(\man) $ is simply identified with $ \man $. A key property of conformal embeddings is that they preserve orthogonality relations at each point $ p \in \man $, since \eref{eq:conformal-relation} clearly implies that $ \tilde{g}_p(u,v) = 0 \iff g_p(u,v) = 0 $ for any $ u,v \in T_p\man $. In particular, for Lorentzian manifolds this means that null vectors for $ \tilde{g} $ coincide with null vectors for $ g $, hence the light cones of $ (\man , g) $ and $ (\man , \tilde{g}) $ are identical: this implies that they have the same casual structure.

The same, however, is not true for geodesics, as elucidated with the formalism of \lref{lemma:difference-tensor}. Denoting as $ \na $ and $ \tilde{\na} $ the Levi--Civita connections associated to $ g $ and $ \tilde{g} $ respectively, they determine a smooth tensor field $ \cont \in \tns{1}{2}(\man) $ which acts on $ X,Y \in \xm $ as:
\begin{equation}
  \cont(X , Y) \defeq \tilde{\na}_Y - \na_X Y
  \label{eq:connection-tensor}
\end{equation}
The components of $ \cont $ in a local frame on $ \man $ are given by \eref{eq:perturbation-difference-tensor}. From the conformal relation \eref{eq:conformal-relation} $ \na_\alpha \tilde{g}_{\mu \nu} = 2 \Omega g_{\mu \nu} \na_\alpha \Omega $, by the metric compatibility of the Levi--Civita connection, thus:
\begin{equation}
  C^\alpha_{\mu \nu} = 2 \delta^\alpha_{(\mu} \na_{\nu)} \log \Omega - g_{\mu \nu} g^{\alpha \beta} \na_\beta \log \Omega
  \label{eq:conformal-difference-tensor}
\end{equation}

\begin{proposition}{Conformal preservation of null geodesics}{}
  Let $ (\man,g) $ and $ (\man , \tilde{g}) $ be two Lorentzian manifolds with conformally-related metrics. Then, a smooth curve $ \gamma $ in $ \man $ is a null geodesic for $ \tilde{g} $ if and only if it is a null geodesic for $ g $.
\end{proposition}

\begin{proofbox}
  \begin{proof}
    Given a geodesic $ \gamma $ for $ g $, its tangent vector field $ X $ satisfied $ \na_X X = 0 $. By \eref{eq:connection-tensor}, then:
    \begin{equation*}
      \tilde{\na}_X X = \na_X X - \cont(X,X) = \cont(X,X)
    \end{equation*}
    which in components becomes:
    \begin{equation*}
      (\na_X X)^\alpha = 2 X^\alpha \na_X \log \Omega - \braket{X,X}_g g^{\alpha \beta} \na_\beta \log \Omega
    \end{equation*}
    Therefore, in general geodesics for $ g $ are \emph{not} geodesics for $ \tilde{g} $, and vice versa. However, if $ \gamma $ is a null geodesic $ \braket{X , X}_g = 0 $, hence the equation above becomes:
    \begin{equation}
      \na_X X = (2 \na_X \log \Omega) X
    \end{equation}
    which is just the equation for a non-affinely-parametrized geodesic. It can be shown that the affine parameter $ \tilde{\lambda} $ of $ \gamma $ for $ \tilde{g} $ is related to the affine parameter $ \lambda $ of $ \gamma $ for $ g $ by:
    \begin{equation*}
      \frac{\dd \tilde{\lambda}}{\dd \lambda} = c \Omega^2
    \end{equation*}
    for a constant $ c \in \R $.
  \end{proof}
\end{proofbox}

Having shown that conformal embeddings preserve the causal structure, turn now to the study of curvature.

\begin{lemma}{Conformal transformation of the curvature tensor}{}
  Let $ (\man,g) $ and $ (\man , \tilde{g}) $ be two Lorentzian manifolds with conformally-related metrics. Then, the components $ \tensor{\tilde{\rr}}{_{\mu \nu \alpha}^\beta} $ and $ \tensor{\rr}{_{\mu \nu \alpha}^\beta} $ of the Riemann tensor for $ \tilde{g} $ and $ g $, respectively, are related by:
  \begin{multline}
    \tensor{\tilde{\rr}}{^\beta_{\mu \nu \alpha}} = \tensor{\rr}{^\beta_{\mu \nu \alpha}} + 2 \delta^\beta_{[\mu} \na_{\nu)} \na_\alpha \log \Omega - 2 g^{\beta \rho} g_{\alpha [\mu} \na_{\nu]} \na_\rho \log \Omega + 2 (\na_{[\mu} \log \Omega) \delta^\beta_{\nu]} \na_\alpha \log \Omega + \\ - 2 (\na_{[\mu} \log \Omega) g_{\nu]\alpha} g^{\beta \rho} \na_\rho \log \Omega - 2 g_{\alpha [\mu} \delta^{\beta}_{\nu]} g^{\rho \sigma} \na_\rho \log \Omega \na_\sigma \log \Omega
    \label{eq:conformal-riemann}
  \end{multline}
  where $ \na $ is the Levi--Civita connection associated to $ g $.
\end{lemma}

Contracting \ref{eq:conformal-riemann}, the conformal transformation law of the Ricci scalar curvature is found:
\begin{equation}
  \tilde{\rr} = \Omega^{-2} \left[ \rr - 2 (d-1) g^{\mu \nu} \na_\mu \na_\nu \log \Omega - (d-2)(d-1) g^{\mu \nu} \na_\mu \log \Omega \na_\nu \log \Omega \right]
  \label{eq:conforma-ricci-scalar}
\end{equation}
where $ \tilde{\rr} \equiv \tilde{g}^{\mu \nu} \tensor{\tilde{\rr}}{^\alpha_{\mu \alpha \nu}} $ and $ \rr \equiv g^{\mu \nu} \tensor{\rr}{^\alpha_{\mu \alpha \nu}} $. It is convenient to define the ``trace-free part" of the Riemann tensor by subtracting all its possible contractions, yielding the \bctxt{Weyl tensor} $ \weyl \in \tns{1}{3}(\man) $, which has components:
\begin{equation}
  \ww_{\mu \nu \alpha \beta} \defeq \rr_{\mu \nu \alpha \beta} + \frac{2}{d - 2} \left( g_{\mu [\alpha} \rr_{\beta] \nu} - g_{\nu [\alpha} \rr_{\beta] \mu} \right) - \frac{2}{(d-1)(d-2)} g_{\mu [\alpha} g_{\beta]\nu} \rr
\end{equation}
This expression if only defined for $ d \geq 3 $, as for $ d = 1,2 $ the Riemann tensor is trivial (recall \eref{eq:riemann-deg-freedom}): in $ d = 2 $ it has a single independent component, given by the Ricci scalar, while in $ d = 1 $ it vanishes identically.
The Weyl tensor has the same symmetries \eref{eq:riemann-symm} of the Riemann tensor, with the additional property of being trace-free in any pair of its indices. Moreover, a crucial property of $ \weyl $ is that it is conformally invariant, as can be shown from \eref{eq:conformal-riemann}:
\begin{equation*}
  \tensor{\tilde{\ww}}{^\beta_{\mu \nu \alpha}} = \tensor{\ww}{^\beta_{\mu \nu \alpha}}
\end{equation*}
For this reason, $ \weyl $ is often called the \emph{conformal (curvature) tensor}. Note that, since indices are naturally raided and lowered with different metrics, the conformal invariance relation depends on the position of the indices: e.g. $ \tilde{\ww}_{\beta \mu \nu \alpha} \equiv \tilde{g}_{\beta \rho} \tensor{\tilde{C}}{^\beta_{\mu \nu \alpha}} = \Omega^2 g_{\beta \rho} \tensor{\ww}{^\beta_{\mu \nu \alpha}} \equiv \Omega^2 \ww_{\beta \mu \nu \alpha} $. The Weyl tensor is then linked to tidal effects, which are indeed conformally invariant, since they remain unchanged when rescaling the volume.

\begin{observation}{Petrov classification}{}
  In the absence of matter, the Einstein field equations nullify the Ricci tensor, hence the only independent components of the Riemann tensor in vacuum are given by the Weyl tensor. For this reason the Weyl tensor is used to classify the vacuum solutions of the Einstein field equations: in particular, in $ d = 4 $ it is possible to take advantage of the symmetry of Hodge operator $ \star : \df{2} \ra \df{2} $, obtaining what is known as the \bcobs{Petrov classification} of vacuum solutions.
% TODO: complete observation on petrov classification
\end{observation}

\subsubsection{Conformal completion of Minkowski spacetime}

Before defining conformal completions of arbitrary spacetimes, it is instructive to first build intuition with the special case of $ d = 4 $ Minkowski spacetime $ (\mink,\eta) $.

To study the ``far-away" behaviour of Minkowski spacetime, it is convenient to work in spherical coordinates $ (t , r , \vartheta , \varphi) $, which are related to the global Cartesian coordinates by:
\begin{equation*}
  x = r \sin \vartheta \cos \varphi
  \quad
  y = r \sin \vartheta \sin \varphi
  \quad
  z = r \cos \vartheta
\end{equation*}
These coordinates span $ \R \times \R^+ \times (0 , \pi) \times (0 , 2\pi) $, but they are not a smooth chart on the whole $ \man $, but only on $ \man - \Pi $, where $ \Pi $ is the closed half-hyperplane defined by $ x \geq 0 $ and $ y = 0 $. In these coordinates, the Lorentz--Minkowski metric takes the form:
\begin{equation*}
  \dd s^2 = - \dd t^2 + \dd r^2 + r^2 \dd \Omega_2^2
\end{equation*}
On the non-angular part $ \R \times \R^+ $, introduce null coordinates $ (u,v) $, respectively the retarded and advanced time, defined by:
\begin{equation*}
  u = t - r
  \qquad
  v = t + r
\end{equation*}
These coordinates span the half-part of $ \R^2 $ defined by $ u < v $, and the metric becomes:
\begin{equation*}
  \dd s^2 = - \dd u \dd v + \frac{1}{4} (v - u)^2 \dd \Omega_2^2
\end{equation*}

In order to perform the conformal completion, consider the finite-range coordinates defined by\footnote{To be precise, note that $ (u,v) $ have dimensions of time, i.e. of length, hence finite-range coordinates should be defined as $ U = \arctan (u / \ell_0) $ and $ V = \arctan (v / \ell_0) $, for some length scale $ \ell_0 $. However, this length scale is completely arbitrary and is eliminated by the conformal transformation, hence it is implicitly set to $ \ell_0 \equiv 1 $.}:
\begin{equation*}
  U = \arctan u
  \qquad
  V = \arctan v
\end{equation*}
The $ \arctan $ function is monotonically increasing, so $ (U,V) $ span the half-part of $ (-\frac{\pi}{2} , \frac{\pi}{2}) \times (-\frac{\pi}{2} , \frac{\pi}{2}) $ defined by $ U < V $. Using standard trigonometric relations, the metric takes the form:
\begin{equation*}
  \dd s^2 = \frac{1}{4 \cos^2 U \cos^2 V} \left[ - 4 \dd U \dd V + \sin^2(V - U) \dd \Omega_2^2 \right]
\end{equation*}
Therefore, the Lorentz--Minkowski metric is conformal to a simpler metric:
\begin{equation}
  \dd \tilde{s}^2 = - 4 \dd U \dd V + \sin^2(V-U) \dd \Omega_2^2
  \label{eq:conf-aux}
\end{equation}
with conformal factor $ \Omega : (-\frac{\pi}{2} , \frac{\pi}{2}) \times (-\frac{\pi}{2} , \frac{\pi}{2}) \ra \R^+ : \Omega(U,V) = 2 \cos U \cos V $. The coordinates $ (U,V) $ are of null type, so consider time-space coordinates $ (\tau , \chi) $ defined by:
\begin{equation*}
  \tau = V + U
  \qquad
  \chi = V - U
\end{equation*}
These range of these coordinates is $ \chi \in (0 , \pi) $ and $ \tau \in (\chi - \pi , \pi - \chi) $, hence drawing the Minkowski spacetime in the $ (\tau , \chi) $ results in a half-diamond shape (see \figref{fig:conformal-minkowski}). Combining all the above coordinate transformations, it is possible to link $ (t,r) $ to $ (\tau , \chi ) $ directly, obtaining:
\begin{equation}
  t = \frac{\sin \tau}{\cos \tau + \cos \tau}
  \qquad \qquad
  r = \frac{\sin \chi}{\cos \tau + \cos \chi}
  \label{eq:esu-embedding}
\end{equation}
Inserting these relations in the Lorentz--Minkowski metric, or directly from \eref{eq:conf-aux}, the conformal metric in the $ (\tau , \chi , \vartheta , \varphi) $ coordinates is found to be:
\begin{equation}
  \dd \tilde{s}^2 = - \dd \tau^2 + \dd \chi^2 + \sin^2 \chi \, \dd \Omega_2^2
  \label{eq:esu-metric}
\end{equation}
The spatial part of the metric is just the standard metric on $ \Sn{3} $ in hyperspherical coordinates $ (\chi , \vartheta , \varphi) $. Moreover, the full metric is regular even upon relaxing the condition on $ \tau $, letting it span the entire $ \R $. Therefore, \eref{eq:esu-metric} is truly a metric $ \tilde{g} $ on $ \esu \defeq \R \times \Sn{3} $, and $ (\esu , \tilde{g}) $ is a spacetimes known as the \bctxt{Einstein static universe} (ESU), which is a solution to the Einstein field equations with $ \Lambda > 0 $ and a cosmological fluid with $ w = 0 $ and $ \rho = \Lambda / (4\pi) $ (see \secref{sssec:matter-content} for details on cosmological fluids). \eref{eq:esu-embedding} are thus a conformal embedding of $ (\mink , \eta) $ into $ (\esu , \tilde{g}) $.

\begin{figure}
  \centering
  \includegraphics[width = 0.30 \textwidth]{conformal-minkowski.png}
  \caption{Conformal diagram of Minkowski spacetime. Curves $ r = \text{const.} $ are drawn in red, while curved $ r = \text{const.} $ are drawn in grey. Figure from \cite{gourgoulhon}.}
  \label{fig:conformal-minkowski}
\end{figure}

Since $ \mink $ and $ \esu $ have the same dimension, $ \mink \in \tp(\esu) $, and its topological closure in $ \esu $ is:
\begin{equation}
  \overline{\mink} = \mink \sqcup \scri^+ \sqcup \scri^- \sqcup \{i^0\} \sqcup \{i^+\} \sqcup \{i^-\}
\end{equation}
as shown in \figref{fig:conformal-minkowski}, where the various types of infinity are:
\begin{itemize}
  \item \bctxt{future null infinity} $ \scri^+ \ssq \esu $, which is the hypersurface defined by $ \tau = \pi - \chi $ and $ \tau \in (0,\pi) $, i.e. $ V = \frac{\pi}{2} $ and $ U \in (-\frac{\pi}{2} , \frac{\pi}{2}) $;
  \item \bctxt{past null infinity} $ \scri^- \ssq \esu $, which is the hypersurface defined by $ \tau = \chi - \pi $ and $ \tau \in (-\pi,0) $, i.e. $ V = -\frac{\pi}{2} $ and $ V \in (-\frac{\pi}{2} , \frac{\pi}{2}) $;
  \item \bctxt{spacelike infinity} $ i^0 \in \esu $, which is defined by $ (\tau , \chi) = (0 , \pi) $, i.e. $ (U,V) = (-\frac{\pi}{2} , \frac{\pi}{2}) $;
  \item \bctxt{future timelike infinity} $ i^+ \in \esu $, which is defined by $ (\tau , \chi) = (\pi , 0) $, i.e. $ (U,V) = (\frac{\pi}{2} , \frac{\pi}{2}) $;
  \item \bctxt{past timelike infinity} $ i^- \in \esu $, which is defined by $ (\tau , \chi) = (-\pi , 0) $, i.e. $ (U,V) = (-\frac{\pi}{2} , -\frac{\pi}{2}) $.
\end{itemize}
Note that, since hyperspherical coordinates $ (\chi , \vartheta , \varphi) $ on $ \Sn{3} $ are singular at $ \chi = 0 , \pi $, the conditions for $ i^0 $ and $ i^\pm $ do indeed define unique points of $ \esu $, whatever the value of $ (\tau , \varphi , \vartheta) $, and not hypersurfaces. Moreover, the vertical left boundary of the half-diamond shape in \figref{fig:conformal-minkowski}, i.e. the segment defined by $ \chi = 0 $ and $ \tau \in (-\pi , \pi) $, is not part of the boundary of $ \mink $, but merely reflects the coordinate singularity at $ \chi = 0 $, which corresponds to the usual coordinate singularity at $ r = 0 $ in ordinary spherical coordinates.

\figref{fig:conformal-minkowski-2d} shows the pathological nature of $ i^0 $: indeed, $ \overline{\mink} $ is self-intersecting at $ i^0 $, so there is no open neighbourhood of $ i^0 $ that is homeomorphic to a neighbourhood in $ \U^4 $, thus it is not possible to endow it with the structure of manifold with boundary (\dref{def:top-man-bound}). On the other hand, $ i^\pm $ are points where the boundary is not smooth, hence $ \mink - \{i^0\} $ can only be defined as a topological manifold with boundary, not a smooth one. For this reason, let $ \scri \equiv \scri^+ \sqcup \scri^- $ and $ \tilde{\man} \defeq \mink \sqcup \scri \ssq \esu $: then, $ \tilde{\man} $ is naturally a smooth manifold with boundary, and its manifold boundary is $ \pa \tilde{\mink} = \scri $.

\begin{figure}
  \centering
  \includegraphics[width = 0.60 \textwidth]{conformal-minkowski-2d.png}
  \caption{Two views of a $ 2 $-dimensional section of $ \esu $ at $ \vartheta = \frac{\pi}{2} $ and $ \varphi = \{0,\pi\} $ (the ``right" and ``left" half of each cylinder, respectively). Again, red curves are $ r = \text{const.} $ curves nad grey curves are $ t = \text{const.} $ ones. Figure from \cite{gourgoulhon}.}
  \label{fig:conformal-minkowski-2d}
\end{figure}

The nature of the various infinities of $ \mink $ is determined by the behaviour of its geodesics: all null geodesics start on $ \scri^- $ and end on $ \scri^+ $, all timelike geodesics start at $ i^- $ and end at $ i^+ $, and all spacelike geodesics start and end at $ i^0 $, as shown in \figref{fig:conformal-minkowski}.

\begin{proposition}{Null conformal boundary of Minkowski spacetime}{}
  Given the spacetime $ (\tilde{\man} , \tilde{g}) $ defined above, $ \scri^+ $ and $ \scri^- $ are null hypersurfaces of $ (\tilde{\man} , \tilde{g}) $.
\end{proposition}

\begin{proofbox}
  \begin{proof}
    Consider first $ \scri^+ $: it is defined by $ V = \frac{\pi}{2} $, so a valid coordinate system on it is $ (U , \vartheta , \varphi) $ with $ U \in (-\frac{\pi}{2} , \frac{\pi}{2}) $. The induced metric $ \tilde{g} \vert_{\scri^+} $ is easily obtained setting $ V = \frac{\pi}{2} $ in \eref{eq:conf-aux}:
    \begin{equation*}
      \dd \tilde{s}^2 \vert_{\scri^+} = \cos^2 U \, \dd \Omega_2^2
    \end{equation*}
    This metric clearly has degenerate signature $ (0 , + , +) $, hence $ \scri^+ $ is a null hypersurface of $ (\tilde{\man} , \tilde{g}) $. Similarly, $ \scri^- $ is defined by $ U = - \frac{\pi}{2} $, so a valid coordinate system on it is $ (V , \vartheta , \varphi) $ with $ V \in (-\frac{\pi}{2} , \frac{\pi}{2}) $, and the induced metric $ \tilde{g}\vert_{\scri^-} $ is:
    \begin{equation*}
      \dd \tilde{s}^2 \vert_{\scri^-} = \cos^2 V \, \dd \Omega_2^2
    \end{equation*}
    This metric has degenerate signature $ (0 , + , +) $, hence $ \scri^- $ is a null hypersurface of $ (\tilde{\man} , \tilde{g}) $.
  \end{proof}
\end{proofbox}

The null character of $ \scri^\pm $ also appears clearly in the conform diagram \figref{fig:conformal-minkowski}.

\subsubsection{Conformal completion of arbitrary spacetimes}

Having elucidated the asymptotic structure of Minkowski spacetime, it is now possible to generalize the construction of conformal completions to arbitrary spacetimes.

\begin{definition}{Conformal completion}{conformal-completion}
  A spacetime $ (\man,g) $ admits a \bcdef{conformal completion} if there exist a Lorentzian manifold with boundary $ (\tilde{\man} , \tilde{g}) $ and a conformal embedding $ \Phi : \man \ra \tilde{\man} $ with $ \Omega \in \cfk{\infty}(\man , \R^+) $ sucht that:
  \begin{enumerate}
    \item \label{item:cc-1} $ \tilde{\man} = \Phi(\man) \sqcup \scri $, with $ \scri \defeq \pa \tilde{\man} $;
    \item \label{item:cc-2} $ \Omega \vert_\scri \equiv 0 $ and $ \dd \Omega \vert_\scri \neq 0 $.
  \end{enumerate}
  $ \scri $ is called the \bcdef{conformal infinity} (or \emph{conformal boundary}) of $ (\man , g) $ within $ (\tilde{\man} , \tilde{g}) $.
\end{definition}

Note that $ \scri $ is the manifold boundary of $ \tilde{\man} $, not the topological boundary of $ \tilde{\man} $ (which is $ \emptyset $, recall \obsref{obs:topological-manifold-boundary}). Moreover, it is not part of the physical spacetime $ \man $, since $ \scri \cap \man = \emptyset $.

Condition \itemref{item:cc-1} expresses that $ \man $ is endowed with some boundary by the conformal embedding $ \Phi $, which in the following is suppressed setting $ \Phi(\man) \equiv \man $. Condition \itemref{item:cc-2} expresses instead that this boundary, which lies ``lies at an infinite distance" with respect to $ g $, is brought to a finite distance with respect to $ \tilde{g} $ by the conformal transformation, which in terms of length elements reads:
\begin{equation*}
  \dd s^2 = \Omega^{-2} \dd \tilde{s}^2
\end{equation*}
The first condition \itemref{item:cc-2} clarifies the ``lies at infinite distance" notion, since $ \Omega^{-2} \ra + \infty $ approaching $ \scri $, while the second condition \itemref{item:cc-2} ensures that $ \scri $ is a regular hypersurface of $ \tilde{\man} $.

In general, unlike the caso of Minkowski spacetime, the conformal boundary $ \scri $ is not necessarily a null hypersurface. Actually, the causal type of $ \scri $ is determined by the cosmological constant.

\begin{proposition}{Causal type of conformal boundaries}{causal-conformal-boundary}
  Let the spacetime $ (\man , g) $ be a solution to the Einstein field equation with a cosmological constant $ \Lambda $, and let it be conformal to a spacetime $ (\tilde{\man} , \tilde{g}) $ with conformal factor $ \Omega \in \cfk{\infty}(\man , \R^+) $ and conformal boundary $ \scri \defeq \pa \tilde{\man} $. If the trace of the EMT is $ T \ra 0 $ as $ \Omega \ra 0 $, then $ \scri $ is a:
  \begin{enumerate}
    \item null hypersurface of $ (\tilde{\man} , \tilde{g}) $ if and only if $ \Lambda = 0 $;
    \item spacelike hypersurface of $ (\tilde{\man} , \tilde{g}) $ if and only if $ \Lambda > 0 $;
    \item timelike hypersurface of $ (\tilde{\man} , \tilde{g}) $ if and only if $ \Lambda < 0 $.
  \end{enumerate}
\end{proposition}

\begin{proofbox}
  \begin{proof}
    From \eref{eq:conforma-ricci-scalar}, inverting the roles of $ \tilde{g} $ and $ g $ using the conformal relation \eref{eq:conformal-relation} and the difference tensor \eref{eq:conformal-difference-tensor}, it is possible to relate the Ricci scalar curvatures $ \tilde{\rr} $ and $ \rr $ of respectively $ \tilde{g} $ and $ g $ as:
    \begin{equation}
      \Omega^2 \tilde{\rr} = \rr - (d-1) \tilde{g}^{\mu \nu} \left[ 2 \Omega \tilde{\na}_\mu \tilde{\na}_\nu \Omega - 8 \pa_\mu \Omega \pa_\nu \Omega \right]
      \label{eq:conformal-ricci}
    \end{equation}
    where $ \tilde{\na} $ is the Levi--Civita connection associated to $ \tilde{g} $. Using \eref{eq:ef-eq-alt} to express $ \rr $, the relation becomes:
    \begin{equation*}
      \Omega^2 \tilde{\rr} = \frac{2}{d-2} \left[ d \Lambda - 8 \pi \gd T \right] - (d-1) \tilde{g}^{\mu \nu} \left[ 2 \Omega \tilde{\na}_\mu \tilde{\na}_\nu \Omega - 8 \pa_\mu \Omega \pa_\nu \Omega \right]
    \end{equation*}
    This equation is, in principle, valid only on $ \man = \tilde{\man} - \scri $. Taking the limit $ \Omega \ra 0 $, with the assumption that $ T \ra 0 $ as $ \Omega \ra 0 $, by continuity it becomes an identity on $ \scri $:
    \begin{equation*}
      \tilde{g}^{\mu \nu} \pa_\mu \Omega \pa_\nu \Omega \vert_\scri = - \frac{d}{4 (d-1) (d-2)} \Lambda
    \end{equation*}
    Now, note that $ \scri $ corresponds to a constant value of $ \Omega $, namely $ \Omega = 0 $, therefore $ \pa_\mu \Omega $ are just the components of the normal vector field $ n \in \Gamma(N\scri) $ to $ \scri $, which is defined as the metric dual of the $ 1 $-form $ \dd \Omega = \pa_\mu \Omega \, \dd x^\mu $ (which is non-vanishing by condition \itemref{item:cc-2} of \dref{def:conformal-completion}). This means that the left-hand side of the above relation is just $ \braket{n,n}_g $, and it constraints the sign of $ \braket{n,n}_g $ to be the opposite of the sign of $ \Lambda $. Given the link between the causal type of a hypersurface and the causal type of its normal vector field, this completes the proof.
  \end{proof}
\end{proofbox}

To generalize the classification of infinities performed in Minkowski spacetime, let $ (\man , g) $ be a time-orientable spacetime admitting a conformal completion $ (\tilde{\man} , \tilde{g}) $ with conformal boundary $ \scri $. Then, it is possible to distinguish two parts of the conformal boundary: the \bctxt{future infinity} of $ (\man , g) $ is the subset $ \scri^+ \ssq \scri $ whose points can be reached from any point in $ \man $ by a future-directed causal curve in $ \tilde{\man} $; while the \bctxt{past infinity} of $ (\man , g) $ is the subset $ \scri^- \ssq \scri $ whose points can be reached from any point in $ \man $ by a past-directed causal curve in $ (\tilde{\man}) $. If $ \scri^+ $ (or $ \scri^- $) is a null hypersurface, then it is called the \emph{future} (or \emph{past}) \emph{null infinity}. If $ \scri^\pm $ are null hypersurfaces and $ \scri = \scri^+ \cup \scri^- $, then $ (\tilde{\man} , \tilde{g}) $ is called a \bctxt{conformal completion at null infinity} of $ (\man , g) $.

\begin{example}{Conformal completion of $ \ads_4 $}{}
  The $ 4 $-dimensional anti-de Sitter spacetime $ \ads_4 $ is $ (\man , g) $, with $ \man \cong \R^4 $. The metric in the so-called \emph{glocal static coordinates} $ (\tau , r , \vartheta , \varphi) $ has components:
  \begin{equation*}
    \dd s^2 = \ell^2 \left[ - (1 + r^2) \dd \tau^2 + \frac{1}{1 + r^2} \dd r^2 + r^2 \dd \Omega_2^2 \right]
  \end{equation*}
  where $ \ell \in \R^+ $ is a constant. Note that $ (\tau , r) \in \R \times \R^+ $ are dimensionless coordinates, while $ (\vartheta , \varphi) \in (0,\pi) \times (0 , 2\pi) $ are the standard spherical coordinates on $ \Sn{2} $. This metric is a solution to the vacuum Einstein field equations with $ \Lambda = - 3 / \ell^2 $, and using the so-called \emph{conformal coordinates} $ (\tau , \chi , \vartheta , \varphi) $, with $ \chi = \arctan r \in (0 , \frac{\pi}{2}) $, it becomes:
  \begin{equation*}
    \dd s^2 = \frac{\ell^2}{\cos^2 \chi} \left[ - \dd \tau^2 + \dd \chi^2 + \sin^2 \chi \, \dd \Omega_2^2 \right]
  \end{equation*}
  Note that $ \chi \in (0,\frac{\pi}{2}) $ only spans one half of the hypersphere $ \Sn{3} $, hence $ \Omega = \ell^{-1} \cos \chi $ determines a conformal completion of $ (\man , g) $ into $ (\tilde{\man} , \tilde{g}) $, where $ \tilde{\man} $ is the part of the ESU $ \esu $ defined by $ \chi \leq \frac{\pi}{2} $, and $ \tilde{g} $ is the standard ESU metric \eref{eq:esu-metric}.

  The conformal boundary $ \scri = \pa \tilde{\man} $ is then the hypersurface $ \chi = \frac{\pi}{2} $, as shown in \figref{fig:ads-conf}: its is spanned by the coordinates $ (\tau , \vartheta , \varphi) $, hence its topology is $ \scri \cong \R \times \Sn{2} $, and it is a timelike hypersurface of $ (\tilde{\man} , \tilde{g}) $, as the induced $ 3 $-metric is clearly Lorentzian (confirming \pref{prop:causal-conformal-boundary}). By the behaviour of radial null geodesics (plotted in \eref{fig:ads-conf}), it is evident that any point of $ \scri $ can be connected to $ \man $ by a future-directed null geodesics as well as by a past-directed one, therefore $ \scri = \scri^+ = \scri^- $, so that $ \scri^\pm $ are timelike hypersurfaces.
\end{example}

\begin{figure}
  \centering
  \includegraphics[width = 0.30 \textwidth]{ads-conf.png}
  \caption{Conformal completion of $ \ads_4 $ spacetime. The conformal boundary $ \scri $ is depicted in yellow, red lines are $ \chi = \text{const.} $ lines (uniformly sampled in terms of $ \tan \chi = \sinh \rho $), green curves are radial null geodeiscs, and the purple curve is a radial timelike geodesic bouncing around $ \chi = 0 $. Figure from \cite{gourgoulhon}.}
  \label{fig:ads-conf}
\end{figure}

\begin{example}{Conformal completion of $ \dss_4 $}{}
  The $ 4 $-dimensional de Sitter spacetime $ \dss_4 $ is $ (\man , g) $, with $ \man \cong \R \times \Sn{3} $. The metric in global coordinates $ (t , \chi , \vartheta , \varphi) $ has components:
  \begin{equation*}
    \dd s^2 = \ell^2 \left[ - \dd t^2 + \cosh^2 t (\dd \chi^2 + \sin^2 \chi \, \dd \Omega_2^2) \right]
  \end{equation*}
  where $ \ell \in \R^+ $ is a constant. Note that $ t \in \R $ is a dimensionless coordinate, while $ (\chi , \vartheta , \varphi) \in (0,\pi)^2 \times (0,2\pi) $ are the standard hyperspherical coordinates on $ \Sn{3} $. This metric is a solution to the vacuum Einstein field equations with $ \Lambda = 3 / \ell^2 $, and using coordinates $ (\tau , \chi , \vartheta , \varphi) $, with $ \tau = 2 \arctan \tanh \frac{t}{2} \in (- \frac{\pi}{2} , \frac{\pi}{2}) $, it becomes:
  \begin{equation*}
    \dd s^2 = \frac{\ell^2}{\cos^2 \tau} \left[ - \dd \tau^2 + \dd \chi^2 + \sin^2 \chi \, \dd \Omega_2^2 \right]
  \end{equation*}
  Defining $ \Omega = \ell^{-1} \cos \tau = (\ell \cosh t)^{-1} $ determines a conformal completion of $ (\man , g) $ into $ (\tilde{\man} , \tilde{g}) $, where $ \tilde{\man} $ is the part of the ESU $ \esu $ defined by $ - \frac{\pi}{2} \leq \tau \leq \frac{\pi}{2} $, and $ \tilde{g} $ is $ \eref{eq:esu-metric} $.

  The conformal boundary $ \scri = \pa \tilde{\man} $, found by $ \Omega \vert_\scri = 0 $, has two connected components: $ \scri^\pm $, which is the hypersurface $ \tau = \pm \frac{\pi}{2} $ of $ \tilde{\man} $. Both are spanned by the coordinates $ (\chi , \vartheta , \varphi) $, hence their topology is $ \scri^\pm \cong \Sn{3} $, and they are both spacelike hypersurfaces of $ (\tilde{\man} , \tilde{g}) $, as the induces $ 3 $-metric is clearly Riemannian (confirming \pref{prop:causal-conformal-boundary}). Finally, it is trivial to see that $ \scri^+ $ and $ \scri^- $ respectively match the definitions of future and past infinity, hence $ (\tilde{\man} , \tilde{g}) $ is a conformal completion at infinity of $ \dss_4 $.
\end{example}

\subsubsection{Asymptotic flatness}
\label{sssec:asymptotic-flatness}

Following the conventions of Penrose \cite{penrose-2}, a spacetime $ (\man , g) $ is \bctxt{asymptotically simple} if it admits a conformal completion such that every null geodesic in $ \man $ has two endpoints in $ \scri $.

This condition, although satisfied by Minkowski, de Sitter and anti-de Sitter spacetimes, is rather restrictive: in particular, it excludes black hole spacetimes, since by definition they contain null geodesics that have no endpoint on $ \scri^+ $. For this reason, define a spacetime $ (\man , g) $ to be \bctxt{weakly asymptotically simple} if there exist $ U \in \tp(\man) $ and an asymptotically simple spacetime $ (\man_0 , g_0) $ with an open neighbourhood $ U_0 \ssq \tilde{\man}_0 $ of $ \scri_0 = \pa \tilde{\man}_0 $ such that $ (U_0 \cap \man_0 , g_0) $ is isometric to $ (U , g) $ (recall \dref{def:riem-isometry}). Of course, for a given weakly asymptotically simple spacetime, there may be different non-overlapping regions $ U $ satisfying this property (e.g. in the maximally extended Kerr spacetime).

\begin{definition}{Asimptotically flat spacetime}{asymptotically-flat}
  A spacetime $ (\man , g) $ is \bcdef{asymptotically flat} if it is weakly asymptotically simple and the Ricci tensor of $ g $ vanishes in an open neighbourhood of $ \scri $.
\end{definition}

By this definition, Minkowski spacetime is asymptotically simple, while de Sitter and anti-de Sitter spacetimes are not.
Penrose \cite{penrose-3} has shown that if $ (\man , g) $ is asymptotically simple and empty, then the Weyl tensor of $ g $ vanishes at $ \scri $, which implies that $ \riem \vert_{\scri} \equiv 0 $: this explains the nomenclature of \emph{asymptotically flat}.

\begin{proposition}{Null boundary of asymptotically flat spacetimes}{}
  The conformal boundary $ \scri $ of an asymptotically flat spacetime $ (\man , g) $ is a null hypersurface of the conformal completion $ (\tilde{\man} , \tilde{g}) $.
\end{proposition}

\begin{proofbox}
  \begin{proof}
    Recall \eref{eq:conformal-ricci}. $ \rr = 0 $ in an open neighbourhood of $ \scri $, by \dref{def:asymptotically-flat}, so the limit $ \Omega \ra 0 $ results into $ \braket{n , n}_g = 0 $: by the same argument of \pref{prop:causal-conformal-boundary}, $ \scri $ is a null hypersurface.
  \end{proof}
\end{proofbox}

\subsection{Formal definition}

\subsubsection{Causal structure}

Let $ (\man , g) $ be a time-orientable spacetime. For any subset $ S \ssq \man $, define:
\begin{itemize}
  \item the \bctxt{chronological future} (or \bctxt{past}) of $ S $ as the subset $ I^\pm(S) \ssq \man $ of points that can be reached from a point of $ S $ by a future-directed (or past-directed) timelike curve of non-zero extent;
  \item the \bctxt{causal future} (or \bctxt{past}) of $ S $ as the subset $ J^\pm(S) \ssq \man $ of points that either are in $ S $ or can be reached from a point of $ S $ by a future-directed (or past-directed) causal curve (recall \dref{def:cauchy-hypersurface}).
\end{itemize}

From these definitions, it is clear that $ S \ssq J^\pm(S) $ and $ I^\pm(S) \ssq J^\pm(S) $. On the contrary, not necessarily $ S \ssq I^\pm(S) $: for instance, if $ \man $ does not contain any closed timelike curve, then $ S \cap I^\pm(S) = \emptyset $ for $ S = \{p\} $, with $ p \in \man $.

\begin{lemma}{Topological properties of causal structures}{tpcs}
  Let $ (\man , g) $ be a time-orientable spacetime, and let $ S \ssq \man $. Then:
  \begin{enumerate}
    \item \label{item:tpcs-1} $ I^\pm(S) \in \tp(\man) $, while $ J^\pm(S) $ is not necessarily closed;
    \item \label{item:tpcs-2} $ \Int J^\pm(S) = I^\pm(S) $;
    \item \label{item:tpcs-3} $ \overline{J^\pm(S)} = \overline{I^\pm(S)} $;
  \end{enumerate}
\end{lemma}

\begin{proofbox}
  \begin{proof}
    See e.g. \S 6.2 of \cite{hawking-hellis} or Chap. 14 of \cite{oneill-1}.
  \end{proof}
\end{proofbox}

By \itemref{item:tpcs-2} and \itemref{item:tpcs-3} of \lref{lemma:tpcs}, it follows that the topological boundaries\footnotemark are $ \pa J^\pm(S) = \pa I^\pm(S) $. A spacetime $ (\man , g) $ is then said to be \bctxt{causally simple} if $ J^\pm(K) \ssq \man $ are closed for every compact subset $ K \ssq \man $.

\footnotetext{Recall that, given a topological space $ \mathcal{X} $, the \emph{topological boundary} of $ S \ssq \mathcal{X} $ is $ \pa S \defeq \overline{S} - \Int S $.}

If the spacetime $ (\man , g) $ admits a conformal completion $ (\tilde{\man} , \tilde{g}) $ with conformal infinity $ \scri $, then, by the above definitions, the future and past infinity of $ \man $ are expressible as:
\begin{equation}
  \scri^\pm = \scri \cap J^\pm(\man)
\end{equation}
Note that $ J^\pm(\man) $ are the causal future and causal past of $ \man $ considered as a subset of $ \tilde{\man} $.

\subsubsection{Black holes}

Consider a spacetime $ (\man , g) $ which admits a conformal completion with conformal boundary $ \scri $: then, the future infinity $ \scri^+ $ is \bctxt{complete} if it is a null hypersurface such that all its generators are complete null geodesics\footnote{A geodesic is \emph{complete} if its affine parameters range on the whole $ \R $. In particular, complete geodesics are indextendible.}. The neighbourhood of $ \scri^+ $ in $ \tilde{\man} $ can be considered as the infinitely-far region reached by outgoing null geodesics: hence, if a null geodesic does not reach this regione, it can be considered as trapped somewhere else in spacetime, i.e. in a black hole regione.

\begin{definition}{Black hole}{black-hole}
  Let $ (\man , g) $ be a time-oriented spacetime with a conformal completion such that its future infinity $ \scri^+ $ is complete. A \bctxt{black hole} is the set of points of $ \man $ that are not in the causal past of $ \scri^+ $:
  \begin{equation}
    \bh \defeq \man - (J^-(\scri^+) \cap \man)
    \label{eq:bh-def}
  \end{equation}
  The topological boundary $ \eh^+ \defeq \pa \bh $ is called the \bctxt{future event horizon} (or simply \emph{event horizon}, if no ambiguity arises).
\end{definition}

From \eref{eq:bh-def}, a black hole is the set of points of $ \man $ from which no future-directed causal curve in $ \tilde{\man} $ reaches $ \scri^+ $. A useful expression for the future event horizon can be found recalling the topological identity $ \pa \bh = \overline{\bh} \cap \overline{\man - \bh} $:
\begin{equation*}
  \eh^+ = \overline{\man - (J^-(\scri^+) \cap \man)} \cap \overline{(J^-(\scri^+) \cap \man)} = \pa (J^{-1}(\scri^+) \cap \man)
\end{equation*}
Now, the boundary of $ J^-(\scri^+) $ in $ \tilde{\man} $ is $ \pa J^-(\scri^+) = \pa (J^-(\scri^+) \cap \man) \cup \scri^+ $, so $ \pa J^-(\scri^+) \cap \man = \pa(J^-(\scri^+) \cap \man) $ and:
\begin{equation}
  \eh^+ = \pa J^-(\scri^+) \cap \man
  \label{eq:event-horizon}
\end{equation}
This means that the future event horizon is the part of the boundary of the causal past of the future infinity $ \scri^+ $ that lies in $ \man $, as illustrated in \figref{fig:bh-def}. From \lref{lemma:tpcs}, $ \eh^+ = \pa I^-(\scri^+) \cap \man $.

\begin{figure}
  \centering
  \includegraphics[width = 0.70 \textwidth]{bh-def.png}
  \hspace{30pt}
  \includegraphics[width = 0.20 \textwidth]{bh-mink.png}
  \caption{On the left, a black hole $ \bh $ defined as the complement of the causal past of the future infinity, $ J^-(\scri^+) $. On the righ, a spurious black hole $ \bh $ in Minkowski spacetime resulting from a conformal completion with an incomplete $ \scri^+ $. Figures from \cite{gourgoulhon}.}
  \label{fig:bh-def}
\end{figure}

\begin{example}{No black hole in Minkowski spacetime}{}
  The Minkowski spacetime $ (\mink , \eta) $ contains no black hole, since all future-directed null geodesics terminate in $ \scri^+ $: this implies that $ \man \ssq J^-(\scri^+) $, so that $ \bh = \emptyset $. More generally, any asymptotically simple spacetime contains no black hole, as evident from definition (see \secref{sssec:asymptotic-flatness}).
\end{example}

If one relaxes the assumption of $ \scri^+ $ being complete in \dref{def:black-hole}, unphysical or ``spurious" black holes emerge. For instance, consider the conformal completion of Minkowski spacetime $ (\man , g) $ from its embedding $ \Phi : \man \ra \esu $ into the ESU (\eref{eq:esu-embedding}), but define now $ \scri^+ $ as the hypersurface of $ \esu $ geiven by $ \tau = \pi - \chi $ and $ \tau \in (0,\frac{\pi}{2}) $, instead of $ \tau \in (0,\pi) $, while keeping $ \scri^- $ unchanged: the manifold with boundary $ \tilde{\man} \equiv \man \cup \scri^+ \cup \scri^- $, equipped with the metric \eref{eq:esu-metric}, is then a conformal completion at null infinity of $ (\man , g) $. With such a $ \scri^+ $, the black hole \eref{eq:bh-def} is non-empty, as shown in \figref{fig:bh-def}.

\begin{example}{Schwarzschild black hole}{schw-bh}
  The prototype of a black hole is the Schwarzschild black hole, which is studied in \secref{sec:bh-schwarzschild}. In particular, the Schwarzschild spacetime contains a region $ \bh $ satisfying \dref{def:black-hole}, and it can be studied using a conformal diagram of the maximally-extended Schwarzschild spacetime, also known as a \emph{Carter--Penrose diagram} (see \secref{sssec:bh-schwarzshild-conformal-diag}), which is pictured in \figref{fig:schw-cart-pen}. This diagram offers a compactified view of the entire spacetime, with each point being a $ 2 $-sphere of metric area $ 4 \pi r^2 $.

  The maximally-extended Schwarzschild spacetime has two causally-disconnected regions where $ 2m < r < + \infty $, denoted by $ \man_\text{I} $ and $ \man_\text{III} $ in \figref{fig:schw-cart-pen}: both are asymptotically flat, and their large-$ r $ parts resemble the conformal diagram for Minkowski spacetime (\eref{fig:conformal-minkowski}). Their future conformal boundaries $ \scri^+_\text{I} $ and $ \scri^+_\text{III} $ are null hypersurfaces of $ (\tilde{\man} , \tilde{g}) $, confirming \pref{prop:causal-conformal-boundary} since the Schwarzschild solution has $ \Lambda = 0 $: the full future conformal boundary $ \scri^+ = \scri^+_\text{I} \cup \scri^+_\text{III} $ is thus a null hypersurface with two connected components.

  From \figref{fig:schw-cart-pen}, it is clear that $ J^-(\scri^+) \cap \man $ does not cover the whole $ \man $, hence there is a black hole denoted by $ \bh $. Its event horizon $ \eh^+ $ is a $ r = 2m $ hypersurface, where $ m $ is the mass parameter of the Schwarzschild solution.
\end{example}

\begin{figure}
  \centering
  \includegraphics[width = 0.70 \textwidth]{bh-schw.png}
  \caption{Maximally-extended Schwarzschild spacetime, depicted in a Carter--Penrose diagam. The yellow area is the causal past of the future null infinity $ J^-(\scri^+) $, divided into three parts by $ r = 2m $ hypersurfaces, the grey area is the black hole $ \bh $, green lines are radial null geodesics and orange zigzag lines at $ r = 0 $ denote curvature singularities. Figure from \cite{gourgoulhon}.}
  \label{fig:schw-cart-pen}
\end{figure}

In general, the black hole region in a given spacetime $ (\man , g) $ with a disconnected future infinity $ \scri^+ $ may depend on the choice of connected components of $ \scri^+ $, even if each component is assumed complete. For instance, consider the maximally-extended Schwarzschild spacetime: the regions $ \man_\text{I} $ and $ \man_\text{III} $, in which $ r \in (2m , +\infty) $, are both asymptotically flat and causally disconnected, since no causal curve can connect a point in $ \man_\text{I} $ to one in $ \man_\text{III} $ or vice versa. Then, choosing WLOG an observer in $ \man_\text{I} $, it is possible to consider a conformal completion $ (\tilde{\man} , \tilde{g}) $ of $ (\man , g) $ such that $ \scri = \scri_\text{I} $: the corresponding black hole region $ \bh $ is then larger from that in \exref{ex:schw-bh}, since it now contains $ \man_\text{III} $ too, as shown in \figref{fig:bh-schw-alt}. From the point of view of the observer in $ \man_\text{I} $, these two choices are physically indistinguishable, since no particle can travel from $ \man_\text{III} $ to $ \man_\text{I} $ anyway.

\begin{figure}
  \centering
  \includegraphics[width = 0.70 \textwidth]{bh-schw-alt.png}
  \caption{Maximally-extended Schwarzschild spacetime, but with a different conformal completion, so that $ \scri^+ $ lies only at the boundary of the asymptotically-flat region $ \man_\text{I} $. Figure from \cite{gourgoulhon}.}
  \label{fig:bh-schw-alt}
\end{figure}

\subsubsection{White holes}

Inverting past and future in the definition of a black hole \eref{eq:bh-def}, one defines the so-called \bctxt{white hole} of a spacetime $ (\man , g) $ with a conformal completion with complete past infinity $ \scri^- $:
\begin{equation}
  \wh \defeq\man - (J^+(\scri^-) \cap \man)
\end{equation}
The white hole region is thus the set of points of $ \man $ that cannot be reached by a future-directed causal curve amanating from $ \scri^- $. The boundary of this white hole is called the \bctxt{past event horizon}:
\begin{equation}
  \eh^- \defeq \pa \wh = \pa J^+(\scri^-) \cap \man
\end{equation}
Carter \cite{carter} introduced the concept of \bctxt{domain of outer comunications} as the part of $ \man $ that lies neither in a black hole nor in a white hole:
\begin{equation}
  \doc{\man} \defeq \man - (\bh \cup \wh) = (J^-(\scri^+) \cap J^+(\scri^-)) \cap \man
\end{equation}
The last equality shows that the domain of outer communications is the set of points from which it is possible to exchange signals with arbitrarily fat regions. Moreover, from the definitions of black and white hole, it follows trivially that:
\begin{equation}
  \pa \doc{\man} = \eh^+ \cup \eh^-
\end{equation}

\begin{example}{Schwarzschild white hole}{}
  In the maximally-extended Schwarzschild spacetime depicted in \figref{fig:schw-cart-pen}, the region $ \wh $ is a white hole, with the past event horizon $ \eh^- $ being a null hypersurface $ r = 2m $.
\end{example}

\subsubsection{Achronal boundaries}

% TODO: add proofs from both gourgoulhon and mtw

From the general definition of black hole \dref{def:black-hole}, it is possible to derive some geometric properties of its boundary $ \eh^+ $, independently of the gravitational theory being used.

Let $ (\man , g) $ be a time-orientable spacetime. Following the conventions of \cite{hawking-hellis}, a set of the type $ \pa J^\pm(S) $, with $ S \ssq \man $, is called an \bctxt{achronal boundary} in $ \man $. In particular, event horizons are achronal boundaries.

\begin{proposition}{Achronal character of achronal boundaries}{achronal-boundary}
  Let $ (\man , g) $ be a time-orientable spacetime. Any achronal boundary in $ \man $ is an \bcprop{achronal set}, i.e. no pair of its points can be connected by a timelike curve of $ \man $.
\end{proposition}

This result justifies the used nomenclature. Note that, in this notion of achronal set, it is not required that the timelike curve lies entirely in the set. Accordingly, an analogous statement is: no timelike curve of $ \man $ meets an achronal boundary at more than one point.

Since $ \eh^+ = \pa \bh $, where $ \dim_\R \bh = d $, one would expect the future event horizon to be an hypersurface of spacetime. However, $ \eh^+ $ is not necessarily a smooth manifold: there may exist some points where $ \eh^+ $ is not differentiable.

\begin{proposition}{Topological nature of achronal boundaries}{achronal-topology}
  Let $ (\man , g) $ be a time-orientable spacetime. Any achronal boundary in $ \man $ is a codimension-$ 1 $ topological submanifold of $ \man $.
\end{proposition}

To be precise, any achronal boundary in $ \man $ is a \emph{Lipschitz submanifold} of $ \man $, i.e. an intermediate between a topological submanifold, i.e. a submanifold of class $ \cfk{0} $, and a differentiable submanifold of class $ \cfk{1} $.

Besides being an achronal codimension-$ 1 $ topological submanifold, the main property of achronal boundaries has been proved by Penrose \cite{penrose-4}.

\begin{theorem}{Achronal boundaries ruled by null geodesics}{achronal-null-geo}
  Let $ (\man , g) $ be a time-orientable spacetime, and let $ \ab \ssq \man $ be a future achronal boundary. Then, $ \ab $ is ruled\footnote{Analogy with a \emph{ruled surface}, i.e. a surface $ \rsscript{S} $ in Euclidean space such that there is a straight line lying in $ \rsscript{S} $ through each point of $ \rsscript{S} $ (e.g. a cylinder or a cone). The analogy holds because straight lines are the geodesics of Euclidean space.} by a family of null geodesics, called the \bcth{generators} of $ \ab $, that:
  \begin{enumerate}
    \item \label{item:fehrng-1} either lie entirely in $ \ab $ or never leave $ \ab $ when followed into the future from the point where they enter $ \ab $;
    \item \label{item:fehrng-2} have no endpoint in the future.
  \end{enumerate}
  Moreover, there is exactly one generator through each point of $ \ab $, except at points where null geodesics enter $ \ab $, called \bcth{crossovers}. A crossover is called \emph{caustic} if it is a point where neighbourhing null geodesics focus and converge while arriving in $ \ab $.
\end{theorem}

The same result holds exchanging future $ \lra $ past.
As a consequence, once a null geodesic has merged with $ \ab $ (at a point where it may intersect other null geodesics), it will remain forever on $ \ab $ and will never intersect any other generator. The set of all crossovers is called the \bctxt{crease set} of $ \ab $.

% TODO: should clarify this concepts, see Chap. 34 of MTW

\begin{proposition}{Achronal boundaries as null hypersurfaces}{achronal-null}
  Let $ (\man , g) $ be a time-orientable spacetime, and let $ \ab \ssq \man $ be an achronal boundary. Then, wherever it is smooth, $ \ab $ is a null hypersurface.
\end{proposition}

\begin{proofbox}
  \begin{proof}
    Assume that $ \ab $ is smooth in some subset $ U \in \tp(\man) $. By \pref{prop:achronal-topology}, $ \ab $ is then a smooth hypersurface in $ U $. According to \tref{th:achronal-null-geo}, there is a null geodesic lying in $ \ab $ through any point of $ \ab \cap U $. This implies the existence of null tangent vectors at any point of $ \ab \cap U $, so that $ \ab $ must be either a null or a timelike hypersurface in $ U $. But $ \ab $ is achronal by \pref{prop:achronal-boundary}, therefore it cannot be timelike. Hence, $ \ab $ is a null hypersurface in $ U $.
  \end{proof}
\end{proofbox}

It can be shown that future event horizons are smooth almost everywhere (e.g. the Schwarzschild black hole in \figref{fig:schw-cart-pen}): the only subset where they are not differentiable is the crease set, i.e. the set of points where null geodesics cross each other while entering $ \eh^+ $ to become null generators.

\section{Event horizons}
\label{sec:event-horizons}

Having established with \pref{prop:achronal-null} that the smooth part of the event horizon of a black hole is a null hypersurface, it is convenient to analyze the geometric properties of this type of hypersurfaces, in order to distinguish between true black hole event horizones and other null hypersurfaces, e.g. future light cones.

\subsection{Null hypersurfaces}

Let $ (\man , g) $ be a spacetime. As any hypersurface, a null hypersurface $ \eh $ can be locally considered as regular a level set (\dref{def:level-set}) of a defining map $ u \in \cfk{\infty}(U) $, with $ U \ssq \man : U \cap \eh \neq 0 $. In particular, the smooth function can be chosen so that, for all $ p \in U $:
\begin{equation*}
  p \in \eh \iff u(p) = 0
  \qquad \qquad
  \dd u \vert_\eh \neq 0
\end{equation*}
The second condition ensures that $ \eh $ is not self-intersecting.

\subsubsection{Null geodesic generators}

Let $ \nn \in \Gamma(N \eh) $ be a smooth vector field normal to $ \eh $. Being $ \eh $ a null hypersurface, $ \nn $ is null, i.e. $ \braket{\nn , \nn}_g = 0 $; moreover, WLOG choose it future-directed. As $ \nn^2 = 0 $, there is no natural normalization for null normals, since any rescaling $ \nn \mapsto \nn' = \alpha \nn $ with $ \alpha \in \R^+ $ (to preserver causal orientation) yields an analogous null normal to $ \eh $.

A priori, the null normal vector field $ \nn $ is defined on $ \eh $ only. However, it is useful to consider $ \nn \in \xv(U) $, with $ U \in \tp(\man) : U \supseteq \eh $, so as to allow for the definition of the covariant derivative $ \na \nn $. To do so, consider a foliation of $ \man $, in general in the vicinity of $ \eh $, by a family of null hypersurfaces: WLOG this foliation can be seen as a family of level set of the above-defined $ u \in \cfk{\infty}(U) $, denoted by $ \{\eh_c\}_{c \in \R} : \eh_c = u^{-1}(c) $, so that $ \eh \equiv \eh_0 $. Then, $ \nn $ can be defined on the part of $ \man $ foliated by $ (\eh_c) $ so that $ \forall p \in U \, \exists c \in \R : \nn_p \in N_p\eh_c $, i.e. for each point in the region there is a leaf of the foliation to which $ \ell $ is a null normal.

Obviously, the foliation $ (\eh_c) $ is not unique, but all geometrical quantities defined hereafter do not depend on the chosen foliation once they are evaluated on $ \eh $. Since $ \eh $ is a hypersurface on which $ u $ is constant, then $ \bs{v} \in T_p\man $ is tangent to $ \eh $ if and only if $ \na_{\bs{v}} u = 0 $, i.e. $ \braket{\grad u , \bs{v}}_g = 0 $ since $ u \in \tns{0}{0}(\man) $. This means that $ \grad u $ is a normal vector field on $ \eh $, so, by uniqueness of the normal direction to a hypersurface, it must be collinear to $ \nn $, and it is always possible to choose $ \rho \in \cfk{\infty}(U) $ such that:
\begin{equation}
  \nn = - e^\rho \grad u
  \label{eq:frobenius-aux}
\end{equation}
The strictly-negative sign of the coefficients is always possible to find by a redefinition of $ u $, and it ensures that if $ u $ increases towards the future, then $ \nn $ is future-directed.

\begin{lemma}[before upper = {\tcbtitle}]{Frobenius' identity}{}
  \begin{equation}
    \dd \ell^\flat = \dd \rho \wedge \ell^\flat
    \label{eq:frobenius}
  \end{equation}
\end{lemma}

\begin{proofbox}
  \begin{proof}
    Writing the metric dual of \eref{eq:frobenius-aux} $ \ell^\flat = - e^\rho \dd u $ in a local coordinate frame yields $ \ell_\alpha = - e^\rho \na_\alpha u $. Taking the covariant derivative:
    \begin{equation*}
      \na_\alpha \ell_\beta = - e^\rho \na_\alpha \rho \na_\beta u - e^\rho \na_\alpha \na_\beta u = \na_\alpha \rho \ell_\beta - e^\rho \na_\alpha \na_\beta u
    \end{equation*}
    Antisymmetrizing and using the torsion-free nature of the Levi--Civita connection:
    \begin{equation}
      \na_{[\alpha} \ell_{\beta]} = \na_{[\alpha} \rho \ell_{\beta]}
      \label{eq:frobenius-components}
    \end{equation}
    The left-hand side are the components of $ \dd \ell^\flat $, again by the torsion-free property, while the right-hand side are the components of $ \dd \rho \wedge \ell^\flat $.
  \end{proof}
\end{proofbox}

This identity is linked to the Frobenius theorem, see e.g. Th. B.3.2 of \cite{wald} of Th. C.2 of \cite{straumann}: a $ 1 $-form satisfies \eref{eq:frobenius} if and only if it defines hyperplanes that are integrable in some hypersurface ($ \eh $ in this case). Contracting the component equation \eref{eq:frobenius-components} with $ \nn $ yields:
\begin{equation*}
  \ell^\mu \na_\mu \ell_\alpha - \ell^\mu \na_\mu \ell_\alpha = \ell^\mu \na_\mu \rho \ell_\alpha - \ell^\mu \ell_\mu \na_\alpha \rho = \ell^\mu \na_\mu \rho \ell_\alpha
\end{equation*}
Since $ \nn $ is null, $ \ell^\mu \na_\alpha \ell_\mu = \na_\alpha (\ell^\mu \ell_\mu) - \ell_\mu \na_\alpha \ell^\mu = - \ell_\mu \na_\alpha \ell^\mu $, i.e. $ \ell^\mu \na_\alpha \ell_\mu = 0 $, thus:
\begin{equation}
  \na_{\nn} \nn = \sg \nn
  \qquad \qquad
  \sg \equiv \na_{\nn} \rho
\end{equation}
This is the equation for a non-affinely parametrized geodesic, hence the integral curves of $ \nn $ are null geodesics. $ \sg $ is often called the \emph{non-affinity coefficient}, and it can be set to $ \sg = 0 $ with a suitable\footnotemark rescaling $ \nn \mapsto \nn' = \alpha \nn $. This analysis allows to state a result analogous to \tref{th:achronal-null-geo}.

\footnotetext{Explicitly, if $ \nn' = \alpha \nn $:
\begin{equation*}
  \na_{\nn'} \nn' = \alpha (\na_{\nn} \alpha + \sg \alpha) \nn = 0
  \quad \iff \quad
  \na_{\nn} \log \alpha = - \sg
  \quad \iff \quad
  \frac{\dd}{\dd \lambda} \log \alpha = - \sg
\end{equation*}
where $ \lambda $ is the parameter of $ \nn $. Therefore, setting $ \alpha(\lambda) = \exp \left( - \int_0^\lambda \dd \bar{\lambda} \, \sg(\bar{\lambda}) \right) $ along each integral curve ensures that it is an affinely-parametrized geodesic. Moreover, under a rescaling $ \ell' = \alpha \ell $, the non-affinity coefficient transforms as:
\begin{equation}
  \sg' = \alpha \sg + \na_{\nn} \alpha
\end{equation}}

\begin{proposition}{Geodesic generators of a null hypersurface}{null-geodesic-generators}
  Let $ (\man , g) $ be a spacetime. Any null hypersurface $ \eh \ssq \man $ is ruled by a family of null geodesics, calle \bcprop{null geodesic generators} of $ \eh $, and each $ \nn \in \Gamma(N \eh) $ is tangent to these null geodesics. Moreover, the null geodesic generators of $ \eh $ form a congruence\footnote{Let $ \man $ be a smooth manifold, and let $ \Sigma \ssq \man $ be a smooth submanifold. Then, a \emph{congruence} in $ \Sigma $ is a family of curves in $ \Sigma $ such that through each point in $ \Sigma $ there passes one and exactly one such curve.} in $ \eh $.
\end{proposition}

This result if not trivial: while it is obvious that integral curves of a null normal of $ \eh $ at null curves tangent to $ \eh $, keep in mind that not all null curves are null geodesics. For a smooth achronal boundary, the geodesic generators defined in \pref{prop:null-geodesic-generators} are the same as those defined in \tref{th:achronal-null-geo}.

\subsubsection{Cross-sections}

\begin{definition}{Cross-section of a null hypersurface}{cross-section-null}
  Let $ (\man , g) $ be a $ d $-spacetime, with $ d \geq 3 $, and let $ \eh $ be a null hypersurface. A \bcdef{cross-section} of $ \eh $ is a $ (d-2) $-submanifold $ \crs \ssq \eh $ such that:
  \begin{enumerate}
    \item \label{item:csnh-1} any $ \nn \in \Gamma(N \eh) $ is nowhere tangent to $ \crs $;
    \item \label{item:csnh-2} each null geodesic generator of $ \eh $ intersects $ \crs $ at most once.
  \end{enumerate}
\end{definition}

If each null geodesic generator of $ \eh $ intersects $ \crs $ exactly one, then $ \crs $ is a \bctxt{complete cross-section} (see \figref{fig:cross-section}). In the following, indices relative to a cross-section range from $ 2 $ to $ d -1 $ and are denoted by Latin letters from the beginning of the alphabet.

\begin{figure}
  \centering
  \includegraphics[width = 0.50 \textwidth]{cross-section.png}
  \caption{Two complete cross-section $ \crs $ and $ \crs' $ of a null hypersurface $ \eh $. Green lines represent null geodesic generators. Figure from \cite{gourgoulhon}.}
  \label{fig:cross-section}
\end{figure}

No restriction is placed on the topology of cross-section. However, note that if $ \eh $ is a black hole event horizon, its complete cross-sections are closed manifolds, i.e. compact manifolds without boundary: the simplest example is $ \crs = \Sn{d-2} $, which is the relevant case for $ d = 4 $ (for $ d > 4 $ other compact topologies are possible).

\begin{proposition}{Spacelike character of cross-sections}{}
  Any cross-section of a null hypersurface is spacelike, i.e. all its tangent vectors are spacelike.
\end{proposition}

\begin{proofbox}
  \begin{proof}
    To prove this proposition, first state a useful lemma.

    \begin{lemma}{Tangent vectors to a null hypersurface}{cs-aux}
      Every non-zero tangent vector $ \vv $ to a null hypersurface is either spacelike or null. Moreover, if it is null, $ \vv $ is tangent to a null geodesic generator, i.e. it is normal to the hypersurface.
    \end{lemma}

    Let $ p \in \crs $ and $ \vv \in T_p\man $ be a non-zero tangent vector to $ \crs $. By \lref{lemma:cs-aux}, $ \vv $ is either spacelike to tangent to the null geodesic generator $ \gamma $ going through $ p $, but then $ \gamma $ would be tangent to $ \crs $, which is not allowed by \itemref{item:csnh-1} of \dref{def:cross-section-null}. This implies that $ \vv $ must be spacelike, which proves that $ \crs \ssq \man $ is a spacelike submanifold.
  \end{proof}
\end{proofbox}

Being $ \crs $ spacelike, the metric induced by $ g $ on $ \crs $ is a Riemannian metric: denote it with $ \bar{q} : T \crs \times T \crs \ra \cfk{\infty}(\crs) : \bar{q} \equiv g \vert_\crs $. Then, through the Gram--Schmidt algorithm, it is possible to find at each $ p \in \crs $ an \emph{orthogonal complement} $ T_p^\perp \crs $ of $ T_p \crs $, which is a $ 2 $-dimensional vector space such that:
\begin{equation}
  T_p\man = T_p \crs \oplus T_p^\perp \crs
  \label{eq:crs-dir-sum}
\end{equation}
where each vector in $ T_p^\perp \crs $ is orthogonal to $ \crs $. Since $ \bar{q} $ is Riemannian, the induced metric on $ T_p^\perp \crs $ must be Lorentzian with signature $ (-,+) $, so to ensure that $ g_p = g \vert_{T^\perp \crs} \oplus \bar{q}_p $ has signature $ (-,+,\dots,+) $. Moreover, $ \crs \ssq \eh $ means that a null normal $ \nn $ of $ \eh $ is orthogonal to any vector tangent to $ \crs $, i.e. $ \nn_p \in T_p^\perp \crs \,\, \forall p \in \crs $: one can then supplement $ \nn_p $ with another null vector to form a null basis of $ T_p^\perp \crs $. It can be shown that, at each $ p \in \crs $, there exists a unique such complement $ \kk_p \in T_p\man $ with normalization $ \braket{\kk_p , \nn_p} = -1 $, so denote with $ \kk \in \xm : \kk(p) = \kk_p $ the associated vector field: note that, while $ \nn $ is independent of the choiche of the cross-section $ \crs $ of $ \eh $ through $ p $, $ \kk $ does depend on $ \crs $. Then, the induced metric $ \bar{q} $ on $ \crs $ can be extended to a metric on $ \man $ as:
\begin{equation}
  q \defeq g + \nn^\flat \otimes \kk^\flat + \kk^\flat \otimes \nn^\flat
  \label{eq:crs-metric}
\end{equation}
In a local coordinate frame, this metric has components $ q_{\alpha \beta} = g_{\alpha \beta} + \ell_\alpha k_\beta + k_\beta \ell_\alpha $. Then, at each $ p \in \crs $, given $ \uu , \vv \in T_p\man $:
\begin{equation}
  q(\uu , \vv) = q(\uu^\parallel , \vv^\parallel)
\end{equation}
where $ \uu^\parallel , \vv^\parallel \in T_p \crs $ are the tangent parts of $ \uu $ and $ \vv $ to $ \crs $, which are uniquely defined by the direct sum \eref{eq:crs-dir-sum}. This means that the metric dual of $ q $ is an orthogonal projector onto $ \crs $:
\begin{equation}
  \pi_\crs \defeq \id_\man + \nn \otimes \kk^\flat + \kk \otimes \nn^\flat
\end{equation}
In a local coordinate frame, this projector has components $ \pi^\alpha_\beta = \delta^\alpha_\beta + \ell^\alpha k_\beta + k^\alpha \ell_\beta $, and it acts on $ T_p\man $ as $ \pi_\crs(\vv) = \vv^\parallel $. In particular, $ \pi_\crs(\nn) = \pi_\crs(\kk) = \bs{0} $.

The null basis $ (\nn , \kk) \in T_p^\perp \crs $ can be orthonormalized as follows:
\begin{equation}
  \nb = \frac{1}{2} \nn + \kk
  \qquad \qquad
  \ssb = \frac{1}{2} \nn - \kk
\end{equation}
so that $ \braket{\nb , \nb}_g = -1 $, $ \braket{\ssb , \ssb}_g = 1 $ and $ \braket{\ssb , \nb}_g = 0 $. Substituting into \eref{eq:crs-metric}, the metric becomes:
\begin{equation}
  q = g + \nb^\flat \otimes \nb^\flat - \ssb^\flat \otimes \ssb^\flat
\end{equation}

\subsubsection{Expansion of cross-sections}

Let $ (\man , g) $ be a $ d $-spacetime, with $ d \geq 3 $, let $ \eh \ssq \man $ be a null hypersurface with null normal $ \nn \in \Gamma(N \eh) $, and let $ \crs \ssq \eh $ be a cross-section of $ \eh $. Consider an infinitesimal $ \varepsilon \geq 0 $, and consider a pair of points $ p , p_\varepsilon \in \eh : p \in \crs $ separated by the vector $ \varepsilon \nn $ (clearly $ p_\varepsilon \in \eh $ since $ \nn $ is tangent to $ \eh $ and $ \varepsilon $ is infinitesimal). Repeating this construction for each point in $ \crs $, keeping $ \varepsilon $ fixed, results in a new $ (d-2) $-dimensional surface $ \crs_\varepsilon $, as in \figref{fig:crs-area}: this is called the \emph{Lie dragging} of $ \crs $ along $ \nn $ by $ \varepsilon $. Clearly, $ \crs_0 = \crs $.

\begin{figure}
  \centering
  \includegraphics[width = 0.70 \textwidth]{cross-section-area.png}
  \caption{Lie dragging of $ \crs $ along $ \nn $ by $ \varepsilon $, here pictured with $ d = 3 $. Figure from \cite{gourgoulhon}.}
  \label{fig:crs-area}
\end{figure}

Since $ p_\varepsilon \in \eh \,\,\forall p \in \crs $ for infinitesimal $ \varepsilon $, then $ \crs_\varepsilon \ssq \eh $. Moreover, being the null direction $ \nn $ transverse to $ \crs_\varepsilon $ by construction, it follows that $ \crs_\varepsilon $ is spacelike too (by \lref{lemma:cs-aux}). Then, at each point $ p \in \crs $ it is possible to define the \bctxt{expansion of $ \crs $ along $ \nn $} as the rate of change of the hyper-area $ \delta A $ of a surface element $ \delta S \crs $ around $ p $ (see \figref{fig:crs-area}):
\begin{equation}
  \theta_{(\nn)} \defeq \lim_{\varepsilon \ra 0} \frac{1}{\varepsilon} \frac{\delta A_\varepsilon - \delta A}{\delta A}
  \label{eq:expansion-def}
\end{equation}
where $ \delta A_\varepsilon $ is the hyper-area of the Lie-dragged surface element $ \delta S_\varepsilon \ssq \crs_\varepsilon $. Note that the hyper-area $ \delta A $ is measured with respect to the induced metric $ q $ on $ \crs $, and similarly for $ \delta A_\varepsilon $ with the Riemmannian metric induced by $ g $ on $ \crs_\varepsilon $. For instance, if $ \delta S \ssq \crs $ is defined by some infinitesimal dispacements $ \delta \xx_{(2)} , \dots , \delta \xx_{(d-1)} $:
\begin{equation*}
  \delta A = \volf_\crs(\delta \xx_{(2)} , \dots , \delta \xx_{(d-1)})
  \qquad \qquad
  \volf_\crs \equiv \ip_{\ssb} \ip_{\nb} \volfg
\end{equation*}
which means that $ \delta A = \volfg (\nb , \ssb , \delta \xx_{(2)} , \dots , \delta \xx_{(d-1)}) $. Choosing adapted coordinates to $ \crs $ and $ \nn $, i.e. $ (x^\alpha) = (u , v , x^2 , \dots , x^{d-1}) $ such that $ \eh $ is locally the set $ u = 0 $, $ \crs $ is the set $ (u , v) = (0,0) $ and $ \nn = \pa_v $, then $ \crs_\varepsilon $ is the set $ (u,v) = (0,\varepsilon) $. Choosing the infinitesimal displacements defining $ \delta S $ to be along the coordinate lines of this frame, the hyper-area element becomes:
\begin{equation}
  \delta A_\varepsilon = \sqrt{\mt{q}(\varepsilon)} \, \delta x_{(2)} \dots \delta x_{(d-1)}
  \label{eq:delta-A-e}
\end{equation}
where $ \mt{q}(\varepsilon) $ is the determinant of the representative matrix of the induced metric $ q_\varepsilon $ on $ \crs_\varepsilon $, computed in the chosen coordinate frame. Plugging this into the definition \eref{eq:expansion-def}:
\begin{equation*}
  \theta_{(\nn)} = \frac{1}{2} \frac{\dd}{\dd \varepsilon} \bigg\vert_{\varepsilon = 0} \ln \mt{q}
\end{equation*}
Since $ \varepsilon $ is the value of the coordinate $ v $ associated to $ \nn $, i.e. the coordinate along which the Lie dragging happens, this is just the Lie derivative of $ \ln \mt{q} $ along $ \nn $ (recall \secref{ssec:lie-der}):
\begin{equation}
  \theta_{(\nn)} = \frac{1}{2} \ld_{\nn} \ln \mt{q}
  \label{eq:expansion-lie-der}
\end{equation}
This is a coordinate-independent definition of the expansion of $ \crs $ along $ \nn $, and with further manipulation it can be rewritten in a local coordinate frame on $ \man $ as:
\begin{equation}
  \theta_{(\nn)} = q^{\mu \nu} \na_\mu \ell_\nu
  \label{eq:expansion-comp}
\end{equation}
where $ q^{\mu \nu} $ are the component of the inverse metric of \eref{eq:crs-metric}. Inserting the expression $ q^{\mu \nu} = g^{\mu \nu} = \ell^\mu k_\nu + k^\mu \ell^\nu $ and using the normalizations of $ \nn $ and $ \kk $ (assuming that $ \nn $ remains null on an open neighbourhood of $ \eh $):
\begin{equation}
  \theta_{(\nn)} = \na_\mu \ell^\mu - \sg
  \label{eq:expansion}
\end{equation}
This expression is completely independent\footnotemark of $ \crs $, thus showing that the expansion at a point $ p \in \eh $ does not depend on the choice of cross-section through $ p $: this justifies the notation $ \theta_{(\nn)} $ in place of $ \theta_{(\nn)}(\crs) $, and from now on $ \theta_{(\nn)} $ is called the \emph{expansion of $ \eh $ along $ \nn $}. From \eref{eq:expansion-lie-der}, the expansion rescales as $ \theta_{(\alpha \nn)} = \alpha \theta_{(\nn)} $.

\footnotetext{This can be shown geometrically too. Let $ \crs $ and $ \crs' $ be two distinct cross-sections of $ \eh $ through $ p $, let $ q \in \crs $ be a point close to $ p $ on $ \crs $, and let $ q' \in \crs' $ be the point on $ \crs' $ along the same null geodesic generator as $ q $, so that $ q $ and $ q' $ are separated by $ \varepsilon \nn $. Then, denoting as $ \delta \xx $ and $ \delta \xx' $ the vectors connecting $ p $ to $ q $ and $ q' $, respectively, it is clear that $ \delta \xx' = \delta \xx + \varepsilon \nn $, so that:
\begin{equation*}
  (\delta \xx')^2 = (\delta \xx)^2 + 2 \varepsilon \braket{\delta \xx , \nn}_g + \varepsilon^2 \nn^2 = (\delta \xx)^2
\end{equation*}
since $ \nn $ is null and normal to any tangent vector to $ \eh $, such as $ \delta \ve{x} $. Therefore, the lengths of all segments from $ p $ do not depend on the cross-section chosen, as long as their other ends lie on the same null geodesic generator of $ \eh $.}


\begin{definition}{Deformation rate of a cross-section}{}
  Let $ \crs \ssq \eh $ be a cross-section of a null hypersurface $ \eh $, and let $ \nn \in \Gamma(N \eh) $ be a null normal to $ \eh $. The \bcdef{deformation rate of $ \crs $ along $ \nn $} is defined as:
  \begin{equation}
    \Theta \defeq \frac{1}{2} \pi_\crs \ld_{\nn} q
  \end{equation}
  where the action of the orthogonal projector $ \pi_\crs $ onto $ \crs $ on the bilinear form $ \ld_{\nn} q $ is defined by:
  \begin{equation}
    \pi_\crs \ld_{\nn} q (\uu , \vv) = \ld_{\nn} q(\pi_\crs(\uu) , \pi_\crs(\vv))
    \quad
    \forall \uu , \vv \in T_p\man
  \end{equation}
\end{definition}

Since $ \pi_\crs(\nn) = \pi_\crs(\kk) = \bs{0} $, in components $ \Theta_{\alpha \beta} = \frac{1}{2} \pi^\mu_\alpha \pi^\nu_\beta \ld_{\nn} g_{\mu \nu} $. Recalling \eref{eq:metric-variation}, then:
\begin{equation}
  \Theta_{\alpha \beta} = \pi^\mu_\alpha \pi^\nu_\beta \na_\mu \ell_\nu
\end{equation}
where \eref{eq:frobenius-aux} and $ \pi^\mu_\alpha \ell_\mu = 0 $ have been used. Comparing this to \eref{eq:expansion-comp} it follows that the trace of $ \Theta $ is just the expansion $ \theta_{(\nn)} $:
\begin{equation}
  \theta_{(\nn)} = g^{\mu \nu} \Theta_{\mu \nu} \equiv \tensor{\Theta}{^\mu_\mu}
\end{equation}
The trace-free part of $ \Theta $ is then called the \bctxt{shear tensor of $ \crs $ along $ \nn $}:
\begin{equation}
  \sigma \defeq \Theta - \frac{1}{d - 2} \theta_{(\nn)} q
\end{equation}
where the $ (d-2)^{-1} $ factor is needed to cancel the trace of the induced metric $ q $ on $ \crs $. Note that, like $ q $, both tensor fields $ \Theta $ and $ \sigma $ are tangent to $ \crs $, in the sense that $ \Theta(\vv , \cdot) = \sigma(\vv , \cdot) \equiv 0 \,\,\forall \vv \in T_p^\perp \crs $. Moreover, they share the same rescaling property as $ \theta_{(\nn)} $ under $ \nn \mapsto \nn' = \alpha \nn $: $ \Theta' = \alpha \Theta $ and $ \sigma' = \alpha \sigma $. These scalings show that both the deformation rate and the shear tensor depend on the choice of the null normal $ \nn $ to $ \eh $; contrary to the expansion, they also depend on the choice of the cross-section $ \crs \ssq \eh $.

\subsubsection{Evolution of the expansion}

\begin{theorem}{Null Raychaudhuri--Landau equation}{}
  Let $ (\man , g) $ be a $ d $-spacetime, and let $ \eh \ssq \man $ be a null hypersurface. Then, any $ \nn \in \Gamma(N \eh) $ satisfies:
  \begin{equation}
    \na_{\nn} \theta_{(\nn)} = \sg \theta_{(\nn)} - \frac{1}{d - 2} \theta_{(\nn)}^2 - \sigma_{ab} \sigma^{ab} - \ricc(\nn , \nn)
    \label{eq:raychaudhuri-landau}
  \end{equation}
  where $ \theta_{(\nn)} $ and $ \sigma $ are the expansion and the shear tensor of $ \eh $ along $ \nn $, and $ \sg $ is non-affinity parameter of $ \nn $.
\end{theorem}

Note that, although $ \sigma $ depends on the choice of cross-section, the shear square $ \sigma_{ab} \sigma^{ab} $ does not, since all other quantities in the null Raychaudhuri--Landau\footnotemark equation do not depend on it.
If $ \nn $ is future-directed, the null Raychaudhuri--Landau equation is an evolution equation for $ \theta_{(\nn)} $ along the null generators of $ \eh $.

\footnotetext{Actually, this is a special case of the null Raychauduri--Landau equation in which the vorticity of the vector field $ \nn $ vanishes. The general case considers a generic congruence of null geodesics (i.e. a family of null geodesics such that one and exactly one passaes through each point of $ \man $), and a null vector field tangent to the geodesics of this congruence has a priori a non-vanishing vorticity tensor $ w $, so that the null Raychauduri--Landau equation gets a $ + w_{ab} w^{ab} $ term on the right-hand side, see e.g. Eq. 4.35 of \cite{hawking-hellis}.}

\begin{proposition}{Null convergence condition}{null-convergence-condition}
  Let $ (\man , g) $ be a $ d $-spacetime, and let $ \eh \ssq \man $ be a null hypersurface. If the \bcprop{null convergence condition} holds on $ \eh $, i.e. if $ \ricc(\vv , \vv) \geq 0 $ for any null $ \vv \in T_p\man : p \in \eh $, then the expansion $ \theta_{(\nn)} $ along any null normal $ \nn \in \Gamma(N\eh) $ associated to an affine parametrization of the null geodesic generators of $ \eh $ cannot increas along $ \nn $, i.e.:
  \begin{equation}
    \na_{\nn} \theta_{(\nn)} \leq 0
  \end{equation}
\end{proposition}

\begin{proofbox}
  \begin{proof}
    Choosing $ \nn $ associated to an affine parametrization of the null geodesic generators of $ \eh $ means that $ \sg = 0 $, hence the null Raychauduri--Landau equation becomes:
    \begin{equation*}
      \na_{\nn} \theta_{(\nn)} = - \frac{1}{d - 2} \theta_{(\nn)}^2 - \sigma_{ab} \sigma^{ab} - \ricc(\nn , \nn)
    \end{equation*}
    The first term of the righ-hand side is manifestly non-positive. To show that $ \sigma_{ab} \sigma^{ab} \geq 0 $, consider a $ q $-orthonormal basis of $ T_p \crs $, so that $ q_{ab} = \delta_{ab} $ and $ q^{ab} = \delta^{ab} $ in such a basis; then:
    \begin{equation*}
      \sigma_{ab} \sigma^{ab} = \sum_{a,b = 2}^{n-1} (\sigma_{ab})^2 \geq 0
      \label{eq:convergence}
    \end{equation*}
    The third term is non-positive due to the null convergence condition, hence $ \na_{\nn} \theta_{(\nn)} \leq 0 $.
  \end{proof}
\end{proofbox}

This property justifies the name \emph{convergence condition}: if $ \theta_{(\nn)} < 0 $ in some part of $ \eh $, so that nearby null geodesic generators are converging, then by \eref{eq:convergence} they cannot subsequently diverge.

If the spacetime $ (\man , g) $ is ruled by General Relativity, from \eref{eq:ef-eq-alt} $ \ricc(\nn , \nn) = 8\pi \gd T(\nn , \nn) $, since $ g(\nn , \nn) = 0 $ being $ \nn $ null. Therefore, the null convergence condition is equivalent to the null energy condition (see \secref{ssec:energy-conditions}).

\subsection{Killing horizons}

\subsubsection{Non-expanding event horizons}

It is now possible to distinguish between event horizons and other null hypersurfaces. Indeed, event horizons are characterized as null hypersurface by the ``no-escape" aspect of the naive definition (\secref{sssec:bh-boundaries}): to enforce the ``localized" aspect, instead, it is reasonable to demand that the cross-sections are closed (i.e. compact and without boundary) and with a constant area, i.e. a vanishing expansion, as natural for black holes \emph{at equilibrium}.

\begin{definition}{Non-expanding horizon}{non-expanding-horizon}
  Let $ (\man , g) $ be a $ d $-spacetime. A null hypersurface $ \eh \ssq \man $ is a \bcdef{non-expanding horizon} if:
  \begin{enumerate}
    \item \label{item:neh-1} its topology is $ \eh \cong \R \times \crs $, where $ \crs $ is a closed $ (d-2) $-submanifold;
    \item \label{item:neh-2} $ \theta_{(\nn)} = 0 $ along any null normal $ \nn \in \Gamma(N \eh) $.
  \end{enumerate}
\end{definition}

Note that, given the scaling law $ \theta_{(\alpha \nn)} = \alpha \theta_{(\nn)} $, if the expansion vanishes for some null normal, then it vanishes for all other null normals. For black holes out of equilibrium (\secref{sec:bh-dynamics}), generally $ \theta_{(\nn)} > 0 $.

Given a complete cross-section $ \crs \ssq \eh $ of a non-expanding horizon $ \eh $, its area with respect to $ g $ is (recall \eref{eq:delta-A-e}):
\begin{equation}
  A \defeq \int_\crs \volf_\crs = \int_\crs \dd x^2 \dots \dd x^{d-1} \sqrt{\mt{q}}
\end{equation}
where $ (x^a) = (x^2 , \dots , x^{d-1}) $ is a coordinate frame on $ \crs $ and $ \mt{q} \equiv \det q_{ab} $ is the determinant of the induced Riemannian metric on $ \crs $, computed in this frame. It can be shown that this area is independent of the choice of the complete cross-section $ \crs \ssq \eh $, hence, $ A \equiv A[\eh] $ is called the \bctxt{area} of the non-expanding horizon $ \eh $ (event though it is a hyper-area, being $ \dim_\R \eh = d - 1 \geq 2 $).

\begin{proposition}{Invariance of the cross-section metric}{neh-cross-sect}
  Let $ (\man , g) $ be a spacetime, and let $ \eh \ssq \man $ be a non-expanding horizon. If the null convergence condition holds on $ \eh $, the extended induced metric on any cross-section $ \crs $ is invariant along the geodesic generators of $ \eh $.
\end{proposition}

\begin{proofbox}
  \begin{proof}
    If $ \eh $ is a non-expanding horizon, the null Raychaudhuri--Landau equation \eref{eq:raychaudhuri-landau} becomes:
    \begin{equation*}
      \sigma_{ab} \sigma^{ab} + \ricc(\nn , \nn) = 0
    \end{equation*}
    where $ \ricc $ is the Ricci tensor of the metric $ g $. Assuming the null convergence condition $ \ricc(\nn , \nn) \geq 0 $, this implies $ \sigma_{ab} \sigma^{ab} = 0 $, which is possible only if $ \sigma \equiv 0 $. Since $ \theta_{(\nn)} = 0 $ too, the deformation rate tensor vanishes identically: $ \Theta = \frac{1}{2} \pi_\crs \ld_{\nn} q = 0 $, which is the thesis.
  \end{proof}
\end{proofbox}

Since $ \eh $ is a null hypersurface, the induced metric $ g \vert_\eh $ is degenerate, hence a priori there is no unique Koszul connection on $ \eh $ associated to $ g \vert_\eh $. However, if $ \eh $ is a non-expanding horizon and the null convergence condition holds on $ \eh $, then it can be shown that the Levi--Civita connection $ \na $ associated to $ g $ induces a unique torsion-free connection $ \tilde{\na} : \xv(\eh) \times \xv(\eh) \ra \xv(\eh) $ on $ \eh $ defined by $ \tilde{\na}_{\uu} \vv = \na_{\uu} \vv \,\,\forall \uu , \vv \in \xv(\eh) $. As a consequence, $ (\eh , g \vert_\eh) $ is a \bctxt{totally-geodesic submanifold} of $ (\man , g) $, since any geodesic of $ (\eh , g \vert_\eh) $ is also a geodesic of $ (\man , g) $.

\subsubsection{Killing horizons}

A special kind of non-expanding horizons, which is of fundamental importance in the theory of stationary black holes, is that of Killing horizons with closed-manifold cross-sections.

\begin{definition}{}{}
  Let $ (\man , g) $ be a spacetime, and let $ \xib \in \xm $ be a Killing vector field of $ \man $. A \bcdef{Killing horizon} is a connected null hypersurface $ \eh $ such that, on $ \eh $, $ \xib $ is normal to $ \eh $.
\end{definition}

Since normals to a null hypersurface are tangent to its null geodesic generators, while a Killing vector is tangent to the orbits of the corresponding isometry group\footnotemark, then it is clear that the null geodesic generators of $ \eh $ are orbits of the isometry group of $ (\man , g) $. Moreover, by definition (a normal to a hypersurface cannot vanish), the Killing vector field is non-vanishing and null on the Killin horizon:
\begin{equation*}
  \xib \vert_\eh \neq 0
  \qquad \qquad
  \xib^2 \vert_\eh = 0
\end{equation*}
It will be shown (REF-TO-TH) that in a stationary spacetime, under some mild hypotheses, the connected part of a black hole event horizon must be a Killing horizon.
% WARN: add ref to Hawking's theorem on Killing horizons

\footnotetext{Let $ G $ be a Lie group with a group action $ \Phi : G \times \mathcal{X} \ra \mathcal{X} $ on some set $ \mathcal{X} $: the \emph{orbit} of $ p \in \mathcal{X} $ is the set $ \{\Phi_g(p) : g \in G\} \ssq \mathcal{X} $.}

\begin{example}{Schwarzschild black hole}{schwarzschild-horizon}
  Consider the Schwarzschild spacetime $ (\man , g) $, with $ \man \cong \R^4 $, with the \bcex{ingoing Eddington--Finkelstein coordinates} (\secref{ssec:eddington-finkelstein-coordinates}) $ (t,r,\vartheta,\varphi) \in \R \times \R^+ \times (0,\pi) \times (0,2\pi) $, so that the Schwarzschild metric \ref{eq:schwarz-sol} takes the form:
  \begin{equation*}
    \dd s^2 = - \left( 1 - \frac{2m}{r} \right) \dd t^2 + \frac{4m}{r} \dd t \dd r + \left( 1 + \frac{2m}{r} \right) \dd t^2 + r^2 \dd \Omega_2^2
  \end{equation*}
  Consider $ u \in \cm $ defined in this coordinate frame as:
  \begin{equation}
    u(t,r,\vartheta,\varphi) = \left( 1 - \frac{r}{2m} \right) \exp \frac{r - t}{4m}
  \end{equation}
  The hypersurface $ u = 0 $, which corresponds to the $ 3 $-cylinder $ r = 2m $ (cylinder as its topology is $ \R \times \Sn{2} $), is the \emph{Schwarzschild horizon} $ \eh $, i.e. the event horizon of the Schwarzschild black hole. Then, the components of $ \grad u = g^{\alpha \mu} \pa_\alpha u \pa_\mu $ are computed to be:
  \begin{equation*}
    \grad u = \frac{1}{4m} e^{(r - t)/(4m)} \left[ - \left( 1 + \frac{r}{2m} \right) \frac{\pa}{\pa t} + \left( 1 - \frac{r}{2m} \right) \frac{\pa}{\pa r} \right]
  \end{equation*}
  Recalling \eref{eq:frobenius-aux}, chooses $ \rho \in \cm $ so that $ \ell^0 = 1 $, i.e.:
  \begin{equation*}
    e^\rho = - \frac{1}{\na^0 u}
    \quad \iff \quad
    \rho = \frac{t - r}{4m} - \log \left( 1 + \frac{r}{2m} \right) + \log (4m)
  \end{equation*}
  From \eref{eq:frobenius-aux}, then, the null normal to $ \eh $ is found to be:
  \begin{equation}
    \nn = \frac{\pa}{\pa t} + \frac{r - 2m}{r + 2m} \frac{\pa}{\pa r}
    \label{eq:schw-ll}
  \end{equation}
  which is easily checked to be null, i.e. $ g(\nn , \nn) = 0 $. Since $ \nn \neq \bs{0} $, this proves that all hypersurfaces $ \eh_c $ determined by $ u $, and in particular $ \eh_0 \equiv \eh $, are null. Moreover, on the horizon $ \nn \heq \pa_t $, which clearly is a Killing vector field since no metric component depends on $ t $: therefore, the Schwarzschild horizon is a Killing horizon.
\end{example}

An important feature of Killing horizons is that the deformation rate of any cross-section along any null normal vanishes. Indeed, consider a cross-section $ \crs $ of a Killing horizon $ \eh $, and let the null normal coincide with the Killing vector field on the horizon, i.e. $ \nn \heq \xib $, so that in a neighbourhood of $ \eh $ it can be expressed as $ \nn = \xib + u \cm $, where $ u $ is the defining function of the horizon and $ \vv \in \xm $. Then:
\begin{equation*}
  \Theta_{\alpha \beta} = \pi^\mu_\alpha \pi^\nu_\beta \na_\mu \xi_\nu + u \pi^\mu_\alpha \pi^\nu_\beta \na_\mu v_\nu + \pi^\mu_\alpha \pi^\nu_\beta w_\nu \na_\mu u \heq \pi^\mu_\alpha \pi^\nu_\beta \na_\mu \xi_\nu
\end{equation*}
where $ \pi^\mu_\alpha \na_\mu u \heq 0 $ since $ u $ is costant on $ \eh $, and so on $ \crs \ssq \eh $. By the Killing equation \eref{eq:killing} $ \pi^\mu_\alpha \pi^\nu_\beta \na_\mu \xi_\nu $ is antisymmetric, while $ \Theta_{\alpha \beta} $ is symmetric (since $ \Theta_{\alpha \beta} = \frac{1}{2} \pi^\mu_\alpha \pi^\nu_\beta \ld_{\nn} g_{\mu \nu} $), hence necessarily $ \Theta \equiv 0 $. By the scaling law of $ \Theta $, it follows that $ \Theta \equiv 0 $ along any null normal to $ \eh $.

\begin{proposition}{Killing horizons as non-expanding horizons}{}
  Any Killing horizon with closed-manifold cross-sections is a non-expanding horizon.
\end{proposition}

\begin{proofbox}
  \begin{proof}
    Since the deformation rate of any cross-section along any null normal vanishes, \itemref{item:neh-2} of \dref{def:non-expanding-horizon} is satisfied. \itemref{item:neh-1} is satisfied instead by the requirement of the cross-sections being closed submanifolds.
  \end{proof}
\end{proofbox}

Note that $ \Theta $ vanishes on any Killing horizon, while to get the same result on a generic non-expanding horizon the null convergence condition is needed (\pref{prop:neh-cross-sect}).

\subsubsection{Surface gravity}

Let $ \sg $ be the non-affinity coefficient of the null normal $ \nn \heq \xib $ to a Killing horizon $ \eh $:
\begin{equation*}
  \na_{\xib} \xib \heq \sg \xib
\end{equation*}
Combining the metric dual with the Killing equation \eref{eq:killing} yields $ \xi^\mu \na_\alpha \xi_\mu \heq - \sg \xi_\alpha $, but notice that $ \xi^\mu \na_\alpha \xi_\mu = \frac{1}{2} \na_\alpha (\xi_\mu \xi^\mu) $, so replacing the covariant derivative with the exterior derivative (since $ \xi^\mu \xi_\mu $ is a scalar field):
\begin{equation}
  \dd \braket{\xib , \xib}_g \heq - 2 \sg \xi^\flat
  \label{eq:surface-gravity-form}
\end{equation}
A relation like this was to be expected: if $ u \equiv \braket{\xib , \xib}_g $ is regular in an open neighbourhood of $ \eh $, i.e. if $ \dd u \neq 0 $, then $ \eh $ is the hypersurface $ u = 0 $, since $ \xib $ is null on $ \eh $, so its gradient must be collinear to the normal $ \xib $ to $ \eh $, i.e. $ \dd u \heq \alpha \xi^\flat $.

\begin{lemma}[before upper = {\tcbtitle}]{Non-affinity coefficient}{}
  \begin{equation}
    \sg^2 \heq - \frac{1}{2} \na_\mu \xi_\nu \na^\mu \xi^\nu
    \label{eq:surface-gravity}
  \end{equation}
\end{lemma}

\begin{proofbox}
  \begin{proof}
    Since $ \xib $ is normal to the hypersurface $ \eh $, by the Frobenius theorem there exists a $ 1 $-form $ a $ such that:
    \begin{equation*}
      \dd \xi^\flat \heq a \wedge \xi^\flat
    \end{equation*}
    which in a local coordinate frame becomes:
    \begin{equation*}
      \na_\alpha \xi_\beta - \na_\beta \xi_\alpha \heq a_\alpha \xi_\beta - a_\beta \xi_\alpha
    \end{equation*}
    By the Killing equation \eref{eq:killing} the left-hand side reduces to $ 2 \na_\alpha \xi_\beta $, so contracting with $ \xib $:
    \begin{equation*}
      2 \xi^\mu \na_\mu \xi_\alpha \heq a_\mu \xi^\mu \alpha_\alpha - \xi^\mu \xi_\mu a_\alpha = a_\mu \xi^\mu \xi_\alpha
    \end{equation*}
    By \eref{eq:surface-gravity-form} the left-hand side is just $ 2 \sg \xi_\alpha $, so $ a_\mu \xi^\mu \heq 2 \sg $. Thus:
    \begin{equation*}
      \begin{split}
        4 \na_\mu \xi_\nu \na^\mu \xi^\nu
        & \heq (a_\mu \xi_\nu - a_\nu \xi_\mu) (a^\mu \xi^\nu - a^\nu \xi^\mu) \\
        & \heq a_\mu a^\mu \xi_\nu \xi^\nu - a_\mu \xi^\mu a_\nu \xi^\nu - a_\nu \xi^\nu a_\mu \xi^\mu + a_\nu a^\nu \xi_\mu \xi^\mu \heq - 8 \sg^2
      \end{split}
    \end{equation*}
    which completes the proof.
  \end{proof}
\end{proofbox}

Although \eref{eq:surface-gravity} is an explicit expression of $ \sg $ in terms of $ \xib $, in actual calculations \eref{eq:surface-gravity-form} is usually employed to compute $ \sg $, since it does not involve any covariant derivative.

Since $ \nn \heq \xib $, the vector field $ \nn $ is a symmetry generator on $ \eh $, hence one may expect that $ \sg $ is invariant along $ \nn $: this turns out to be the case. Indeed, taking the Lie derivative along $ \xib $ of \eref{eq:surface-gravity-form} and using \eref{eq:cartan-formula}:
\begin{equation*}
  \begin{split}
    \ld_{\xib} (\dd \braket{\xib , \xib}_g)
    & = \ip_{\xib} \dd^2 \braket{\xib , \xib}_g + \dd \ip_{\xib} \dd \braket{\xib , \xib}_g = 0 \\
    & \heq - 2 \ld_{\xib} (\sg \xi^\flat) = -2 (\ld_{\xib} \sg) \xi^\flat - 2 \sg \ld_{\xib} \xi^\flat = -2 (\ld_{\xib} \sg) \xi^\flat
  \end{split}
\end{equation*}
where $ \braket{\xib , \xib}_g \heq 0 $ being $ \xib $ null on $ \eh $ and $ \ld_{\xib} \xi^\flat = 0 \impliedby \ld_{\xib} \xib = 0 $, as $ \ld_{\xib} g = 0 $ by definition of Killing vector field. Since $ \nn \heq \xi $ and $ \xi^\flat \neq 0 $, this shows that the non-affinity coefficient is constant along the null geodesic generators of the Killing horizon:
\begin{equation}
  \ld_{\nn} \sg = 0
  \label{eq:sg-const}
\end{equation}
Note that, although $ \sg $ is constant along any null geodesic generators, it may in principle vary from one generator to the other. In order to study under which conditions $ \sg $ is constant on the whole horizon, introduce the \bctxt{null dominance condition}, which states that there exists $ f \in \cm $ such that $ \bs{W} \defeq \bs{G}(\nn) - f \nn \in \xm $ is either zero of future-directed (null or timelike) for any future-directed null $ \nn \in \xm $, where $ \bs{G} \in \tns{1}{1}(\man) $ is the metric dual of the Einstein tensor, so that in local coordinate frame:
\begin{equation}
  W^\alpha \defeq - \tensor{G}{^\alpha_\mu} \ell^\mu - f \ell^\alpha
  \label{eq:null-dominance-condition}
\end{equation}
Note that the null dominance condition implies the null convergence condition, since $ \ricc(\nn , \nn) = - \braket{\bs{W} , \nn}_g \geq 0 $ (since $ \bs{W} $ and $ \nn $ are both future-directed).

\begin{theorem}{Zeroth law of black hole dynamics}{0-bhd}
  Let $ (\man , g) $ be a spacetime, and let $ \eh \ssq \man $ be a Killing horizon defined by the Killing vector field $ \xib \in \xm $. Then, if the null dominance condition holds on $ \eh $, the non-affinity coefficient $ \sg $ associated to the null normal coinciding with $ \xib $ on $ \eh $ is constant over $ \eh $.
\end{theorem}

\begin{proofbox}
  \begin{proof}
    The relation $ \na_{\xib} \xib \heq \sg \xib $ is only valid at the horizon, so, to study the constancy of $ \sg $, it is not possible to apply a generic covariant derivative to it. However, note that the rank-$ (0,3) $ tensor with components $ \lct^{\mu \nu \alpha \beta} \xi_\beta $, where $ \lct^{\mu \nu \alpha \beta} \equiv \sqg \, \epsilon^{\mu \nu \alpha \beta} $ is the Levi--Civita tensor, is indeed tangent to the horizon, since $ (\lct^{\mu \nu \alpha \beta} \xi_\mu) \xi_\alpha = 0 $, thus $ \lct^{\mu \nu \alpha \beta} \xi_\beta \na_\alpha $ is a differential operator that can be safely applied to any equation holding on $ \eh $.

    Equivalently, apply $ \xi_{[\alpha} \na_{\beta]} $ to $ \na_{\xib} \xib \heq \sg \xib $, yielding, in a local coordinate frame:
    \begin{equation*}
      \begin{split}
        \xi_\mu \xi_{[\alpha} \na_{\beta]} \sg + \sg \xi_{[\alpha} \na_{\beta]} \xi_\mu
        & \heq \xi_{[\alpha} \na_{\beta]} (\xi^\nu \na_\nu \xi_\mu) \heq (\xi_{[\alpha} \na_{\beta]} \xi^\nu) (\na_\nu \xi_\mu) + \xi^\nu \xi_{[\alpha} \na_{\beta]} \na_\nu \xi_\mu \\
        & \heq (\xi_{[\alpha} \na_{\beta]} \xi^\nu) (\na_\nu \xi_\mu) - \xi^\nu \tensor{\rr}{_{\nu \mu [\beta}^\rho} \xi_{\alpha]} \xi_\rho
      \end{split}
    \end{equation*}
    where \eref{eq:killing-riemann} was used in the last equality. Since $ \xib $ is hypersurface-orthogonal to the horizon, by Frobenius' theorem $ \xi_{[\mu} \na_\nu \xi_{\beta]} \heq 0 $, so using the Killing equation $ \na_\nu \xi_\beta = - \na_\beta \xi_\nu $:
    \begin{equation}
      \xi_\beta \na_\mu \xi_\nu \heq - 2 \xi_{[\mu} \na_{\nu]} \xi_\beta
      \label{eq:sg-aux}
    \end{equation}
    This allows to rewrite the first term on the right-hand side of the above equation as:
    \begin{equation*}
      (\xi_{[\alpha} \na_{\beta]} \xi^\nu) (\na_\nu \xi_\mu) \heq - \frac{1}{2} (\xi^\nu \na_\alpha \xi_\beta) \na_\nu \xi_\mu \heq - \frac{1}{2} \sg \xi_\mu \na_\alpha \xi_\beta \heq \sg \xi_{[\alpha} \na_{\beta]} \xi_\mu
    \end{equation*}
    which cancels the second term on the left-hand side, leaving only:
    \begin{equation*}
      \xi_\mu \xi_{[\alpha} \na_{\beta]} \sg \heq \xi^\nu \tensor{\rr}{_{\mu \nu [\beta}^\rho} \xi_{\alpha]} \xi_\rho
    \end{equation*}
    On the other hand, applying $ \xi_{[\alpha} \na_{\rho]} $ to \eref{eq:sg-aux}:
    \begin{equation*}
      (\xi_{[\alpha} \na_\rho \xi_\beta) \na_\mu \xi_\nu + \xi_\beta \xi_{[\alpha} \na_{\rho]} \na_\mu \xi_\nu \heq -2 (\xi_{[\alpha} \na_{\rho]} \xi_{[\mu}) \na_{\nu]} \xi_\beta - 2 (\xi_{[\alpha} \na_{\rho]} \na_{[\nu} \xi_{|\beta|}) \xi_{\mu]}
    \end{equation*}
    Using \eref{eq:sg-aux} repeatedly, the first term on the left-hand side cancels the first term on the righ-hand side, thus, using \eref{eq:killing-riemann}:
    \begin{equation*}
      - \xi_\beta \tensor{\rr}{_{\mu \nu [\rho}^\sigma} \xi_{\alpha]} \xi_\sigma \heq 2 \xi_{[\mu} \tensor{\rr}{_{\nu] \beta [\rho}^\sigma} \xi_{\alpha]} \xi_\sigma
    \end{equation*}
    Contracting $ \beta $ and $ \rho $ makes the left-hand side vanish, so:
    \begin{equation*}
      - \xi_{[\mu} \tensor{\rr}{_{\nu]}^\sigma} \xi_\sigma \xi_\alpha \heq \xi_{[\mu} \tensor{\rr}{_{\nu]\beta \alpha}^\sigma} \xi^\beta x_\sigma
    \end{equation*}
    The right-hand side of this equation has the same form of the right-hand side of the above equation (since $ \rr_{\mu \nu \alpha \beta} = \rr_{\alpha \nu \mu \beta} $), hence combining them:
    \begin{equation}
      \xi_{[\alpha} \na_{\beta]} \sg \heq - \xi_{[\alpha} \rr_{\beta] \rho} \xi^\rho
    \end{equation}
    The right-hand side vanishes by the null dominance condition, so there only remains $ \xi_{[\alpha} \na_{\beta]} \sg \heq 0 $, which is the thesis.
    % TODO: shoul prove the implictation from the null dominance condition
  \end{proof}
\end{proofbox}

In the context of Killing horizon, in light of \tref{th:0-bhd}, the non-affinity coefficient $ \sg $ is called the horizon's \bctxt{surface gravity}. Note that, if gravity is described by General Relativity, the null dominance condition with $ f = \Lambda $ is equivalent to the \emph{null dominant energy condition}, which states that the vector field with components $ - \tensor{T}{^\mu_\nu} \ell^\nu $ is either zero or future-directed (null or timelike) for any future-directed null $ \nn \in \xm $, where $ T_{\mu \nu} $ is the energy-momentum tensor of matter fields. The null dominant energy condition is implied by the standard dominant energy condition (\secref{ssec:energy-conditions}).

The constancy of $ \sg $ over $ \eh $ (assuming the null dominance condition) allows to classify Killing horizons in two categories:
\begin{enumerate}
  \item if $ \sg = 0 $, then $ \xib $ is a geodesic vector on $ \eh $, and $ \eh $ is called a \bctxt{degenerate Killing horizon};
  \item if $ \sg \neq 0 $, then $ \eh $ is called a \bctxt{non-degenerate Killing horizon}.
\end{enumerate}

\begin{example}{Schwarzschild and Kerr black holes}{}
  For the Schwarzschild black hole, $ \nn $ is given by \eref{eq:schw-ll} (in ingoing Eddington--Finkelstein coordinates), so its covariant derivative is computed to be:
  \begin{equation}
    \na_{\nn} \nn = \frac{4m}{(r + 2m)^2} \frac{\pa}{\pa t} + \frac{4m (r - 2m)}{(r + 2m)^3} \frac{\pa}{\pa r} = \frac{4m}{(r + 2m)^2} \nn
    \label{eq:in-edd-fink}
  \end{equation}
  Therefore, for any leave $ \eh_c $ of the foliation, $ \sg = 4m / (r + 2m)^2 $. In particular, for the event horizon $ \eh $:
  \begin{equation}
    \sg = \frac{1}{4m}
  \end{equation}
  Note that this value agrees with $ \sg = \na_{\nn} \rho $ (see \exref{ex:schwarzschild-horizon}). Clearly, $ \sg $ is a constant, a trivial result given that the Schwarzschild black hole is spherically-symmetric. A general result, valid for a Kerr--Newman black hole with mass parameter $ m $, angular momentum parameter $ a $ and charge parameter $ q $ is (Eq. 12.5.4 of \cite{wald}):
  \begin{equation}
    \sg = \frac{\sqrt{m^2 - a^2 - q^2}}{2m \left( m + \sqrt{m^2 - a^2 - q^2} \right) - q^2}
  \end{equation}
  which is constant too. Note that, for $ a = 0 $ and $ q = 0 $, it reduces to the Schwarzschild case. Moreover, setting $ q = 0 $, for $ a < m $ the Kerr horizon is a non-degerate Killing horizon, while it becomes degenerate for $ a = m $.
\end{example}

To see why $ \sg $ is interpreted as a surface gravity, consider a non-degenerate Killing horizon $ \eh $, so that $ \sg \neq 0 $, and let $ p \in \eh $. Given $ \vv \in T_p\man $ transverse to $ \eh $, i.e. not tangent to $ \eh $, by \eref{eq:surface-gravity-form}:
\begin{equation*}
  \na_{\vv} \braket{\xib , \xib}_g = - 2 \sg \braket{\xib , \vv}_g
\end{equation*}
The right-hand side does not vanish, as $ \sg \neq 0 $ and $ \braket{\xib , \vv}_g \neq 0 $ (being $ \vv $ not tangent to $ \eh $), thus $ \na_{\vv} \braket{\xib , \xib}_g \neq 0 $, the derivative of $ \xib^2 $ along any direction transverse to $ \eh $ does not vanish. Since $ \xib^2 \heq 0 $, this means that, in the vicinity of $ \eh $, on one side $ \xib^2 < 0 $ and on the other side $ \xib^2 > 0 $. Therefore, in the vicinity of a non-degenerate Killing horizon, the associated Killing vector field is timelike on one side, null on the horizon and spacelike on the other side.

Focus on the side of $ \eh $ where $ \xib $ is timelike, and set the ``norm" $ V \equiv \sqrt{- \braket{\xib , \xib}_g} > 0 $. Then, it can be shown that:
\begin{equation}
  \sg^2 = \lim_\eh \na_\mu V \na^\mu V
  \label{eq:sg-lim}
\end{equation}
where $ \lim_\eh $ is the limit as one approaches $ \eh $ from the timelike-$ \xib $ side, which implies $ V \ra 0 $. In this region where $ \xib $ is timelike, define:
\begin{equation}
  \uu \defeq \frac{1}{V} \xib
\end{equation}
which is a future-directed unit timelike vector field (since by convention $ \xib $ is future-directed and null on $ \eh $, otherwise take $ - \xib $). This vector field $ \uu $ can be considered as the $ d $-velocity of an observer whose worldline if a integral curve of $ \xib $, i.e. an orbit of the isometry group generated by $ \xib $. Since spacetime is not changing along this worldline (by the nature of Killing vector fields), this may be called a \emph{stationary observer}. Its $ d $-acceleration then is:
\begin{equation*}
  \aab = \na_{\uu} \uu = \na_{V^{-1} \xib} (V^{-1} \xib) = V^{-1} \left[ -V^{-2} (\na_{\xib} V) \xib + V^{-1} \na_{\xib} \xib \right]
\end{equation*}
Since $ \xib $ is a symmetry generator, $ \na_{\xib} = 0 $, which can be shown as:
\begin{equation*}
  \na_{\xib} V = \xi^\mu \na_\mu \sqrt{- \xi_\nu \xi^\nu} = - \frac{1}{2 \sqrt{- \xi_\nu \xi^\nu}} \xi^\mu \na_\mu (\xi_\nu \xi^\nu) = - \frac{1}{V} \xi^\mu \xi^\nu \na_\mu \xi_\nu = 0
\end{equation*}
since $ \na_{(\mu} \xi_{\nu)} = 0 $ (Killing's equation). Thus, using again the Killing equation:
\begin{equation*}
  a_\mu = \frac{1}{V^2} \xi^\alpha \na_\alpha \xi_\mu = - \frac{1}{V^2} \xi^\mu \na_\alpha \xi_\mu = - \frac{1}{V^2} \na_\alpha (\xi_\mu \xi^\mu) = \frac{1}{2V^2} \na_\alpha V^2 = \na_\alpha \log V
\end{equation*}
Therefore, the $ d $-acceleration can be expressed as a gradient:
\begin{equation}
  \aab = \grad \log V
\end{equation}
Since $ \braket{\aab , \uu}_g = 0 $, the $ d $-accelaration must spacelike, i.e. $ a_\mu a^\mu > 0 $, so it is possible to define its modulus as:
\begin{equation*}
  a \equiv \sqrt{\braket{\aab , \aab}_g} = \sqrt{\na_\mu \log V \na^\mu \log V} = \frac{1}{V} \sqrt{\na_\mu V \na^\mu V}
\end{equation*}
Recalling \eref{eq:sg-lim}, it is then possible to express $ \sg $ in terms of $ a $:
\begin{equation}
  \sg = \lim_\eh V a
\end{equation}
where $ \lim_\eh $ means that the worldline of the observer whose acceleration is $ \aab $ is taken arbitrarily close to the horizon $ \eh $. Since $ V \ra 0 $ approaching $ \eh $, it must be $ a \ra + \infty $ in order to have a finite $ \sg $: this means that the acceleration felt by an observer (i.e. the ``gravity") diverges as it approaches the horizon. In this sense, the physical surface gravity of $ \eh $ is infinite. However, the rescaled acceleration $ V a $ remains finite approaching $ \eh $, and it tends to $ \sg $: for this reason, $ \sg $ is regarded as the \emph{surface gravity} of the Killing horizon $ \eh $.

The distinction between the surface gravity $ \sg $ and the actual gravity $ a $ measured \emph{locally}, i.e. by an observer at rest with respect to $ \eh $ and infinitely close to it. However, in the special case of a Schwarzschild black hole, $ \sg $ can be interpreted as a physical force (per unit mass) measured by a \emph{distant} observer, for which $ \xib $ is timelike in the entire region outside the Killing horizon\footnotemark: in this case, $ \sg $ can be identified with the magnitude of the force exerted by an observer ``at infinity" to hold in place a particle of unit mass close to $ \eh $.

\footnotetext{This is not true for a rotating Kerr black hole, since $ \xib $ becomes null at come ``light-cylinder" outside $ \eh $ and is then spacelike away from it.}

\section{Stationary black holes}
\label{sec:bh-stationary}

Having defined the general notion of black hole in \secref{sec:bh-def}, focus now on the equilibrium case, i.e. black holes in stationary spacetimes.

\subsection{Stationary spacetimes}

\begin{definition}{Stationary spacetime}{stationary-spacetime}
  A spacetime $ (\man , g) $ is called \bcdef{strictly stationary} if it is invariant under the action of $ \R_+ $ and the associated Killing vector field $ \xib \in \xm $ is everywhere timelike.

  If $ (\man , g) $ admits a conformal completion and $ \xib $ is timelike only in the vicinity of the conformal boundary $ \scri $, then the spacetime is called \bcdef{stationary}.
\end{definition}

The condition on the Killing vector fields (generator of $ \R_+ $) is equivalent to requiring that the orbits of the group action are everywhere timelike curves (or only in the vicinity of $ \scri $). The distinction between strictly stationary and stationary spacetimes is due to the fact that, in presence of a black hole, $ \xib $ cannot be timelike everywhere.

A priori, the Killing vector field $ \xib $ is not unique: $ \R_+ $ admits a reparametrization $ t \mapsto \alpha t $, with $ \alpha \in \R - \{0\} $, which corresponds\footnotemark to $ \xib \mapsto \alpha^{-1} \xib $. When the spacetime admits a conformal boundary $ \scri $ (required in presence of a black hole), then $ \xib $ is uniquely fixed by choosing that near $ \scri $ it is future-directed and $ \braket{\xib , \xib}_g \mapsto -1 $: indeed, this fixes $ \alpha = +1 $, in order to preserve the norm and the causal orientation.

If the $ d $-spacetime $ (\man , g) $ is stationary, then it is possible to choose adapted coordinates $ (x^\alpha) = (t , x^1 , \dots , x^{d-1}) $ such that $ \xib = \pa_t $, which means that $ t $ is an ignorable coordinates, i.e. $ \pa_t g_{\alpha \beta} = 0 $.

\footnotetext{Given a connected $ 1 $-dimensional Lie group $ G $, its elements can be labelled as $ g_t \in G : t \in \R $, with $ g_0 = e $. Then, setting the notation for the group action on a smooth manifold $ \man $ as $ \Phi_t \equiv \Phi_{g_t} : \man \ra \man $, the orbit of $ p \in \man $ is a curve $ \gamma_p : \R \ra \man : \gamma_p(t) = \Phi_t(p) $ (which may be degenerate if $ p $ is a fixed point), and the \emph{generator} $ \xib \in \xm $ of $ G $ is the velocity vector field $ \xib_p \defeq \gamma_p'(0) $. Under a rescaling $ t \mapsto \alpha t $, the velocity rescales as $ \gamma_p'(0) \mapsto \alpha^{-1} \gamma_p'(0) $, hence the generator rescales as $ \xib \mapsto \alpha^{-1} \xib $.}

\begin{definition}{Static spacetime}{static-spacetime}
  A spacetime $ (\man , g) $ is called (strictly) \bcdef{static} if it is (strictly) stationary and the Killing vector field $ \xib \in \xm $ generating the stationary group action is orthogonal to a family of hypersurfaces.
\end{definition}

If the $ d $-spacetime $ (\man , g) $ is static, then it is possible to choose adapted coordinates such that $ \xib = \pa_t $ and the hypersurfaces $ t = \text{const.} $ are orthogonal to $ \xib $, i.e. $ \xi^\flat = g(\xib , \xib) \dd t $. This orthogonality condition means that $ g_{0i} = 0 \,\,\forall i \in \{1 , \dots , d-1\} $, hence in these adapted coordinates the metric of the static spacetime reads:
\begin{equation}
  \dd s^2 = g(\xib , \xib) \dd t^2 + g_{ij} \dd x^i \dd x^j
\end{equation}
Clearly, this metric is invariant under $ t \mapsto -t $, which means that a static spacetime is \emph{time-reflection symmetric}.
In loose terms, a spacetime is stationary if ``nothing changes with time", and it is static if in addition ``nothing moves". An example of a stationary spacetime that is not static is a spacetime containing a steadily-rotating body, e.g. a star of a black hole.

\subsubsection{Black holes in stationary spacetimes}

Consider a spacetime $ (\man , g) $ which contains a black hole: in particular, $ (\man , g) $ admits a conformal completion $ (\tilde{\man} , \tilde{g}) $ with a complete future infinity $ \scri^+ $ (\dref{def:black-hole}). Assuming $ (\man , g) $ to be stationary, it is invariant under the isometry group $ \R_+ $ generated by the Killing vector field $ \xib \in \xm $: this can then be extended smoothly to the boundary $ \pa \tilde{\man} = \scri $, thereby defining a group action $ \Phi : \R_+ \times \tilde{\man} \ra \tilde{\man} $ on the whole conformal completion.

\begin{proposition}{Stationary event horizon}{}
  Let $ (\man , g) $ be a stationary spacetime containing a black hole. Then, the black hole event horizon $ \eh $ is globally invariant under the action of the stationary group $ \R_+ $.
\end{proposition}

\begin{proofbox}
  \begin{proof}
    By the above discussion, the stationary group $ \R_+ $ acts on the whole conformal completion $ \tilde{\man} $, so on the conformal infinity $ \scri $ too.
    Let $ t \in \R $ and $ p \in \scri^+ $. Since $ \R_+ $ is an isometry group of $ g $ and $ p $ is infinitely remote from points of $ \man $, then $ \Phi_t(p) $ must be infinitely remote from any point of $ \man $ too: thus, $ \Phi_t(\scri^+) \ssq \scri $. Since $ \Phi_0(\scri^+) = \scri^+ $, by continuity $ \Phi_t(\scri^+) \ssq \scri^+ $, i.e. $ \scri^+ $ is globally invariant under the stationary group action, which means that its causal past $ J^-(\scri^+) $ is invariant too. By \eref{eq:event-horizon}, this concludes the proof.
  \end{proof}
\end{proofbox}

The fact that $ \eh $ is \emph{globally} invariant stresses the fact that it is invariant as a whole, not that each point of $ \eh $ is a fixed point of the group action.

Now, assume that $ \eh $ is smooth (which can be shown to be the case in a stationary context): by \pref{prop:achronal-null}, it is then a null hypersurface. $ \eh $ is globally invariant under the isometry group if and only if its Killing generator $ \xib $ is tangent to $ \eh $, but by \lref{lemma:cs-aux} a timelike vector cannot be tangent to a null hypersurface\footnotemark, thus $ \xib $ is either null or spacelike on $ \eh $. Given a connected component $ \eh_0 \ssq \eh $, if $ \xib $ is null on all $ \eh_0 $, then $ \eh_0 $ is called \bctxt{non-rotating}, while if $ \xib $ is spacelike on some part of $ \eh_0 $, then $ \eh_0 $ is called \bctxt{rotating}.

\footnotetext{This justifies the distinction between stricly stationary and stationary spacetimes, since a black hole spacetime cannot be strictly stationary.}

Intuitively, the event horizon of a black hole in a stationary spacetime should have vanishing expansion along its null normals: indeed, if $ \theta_{(\nn)} \neq 0 $, the area of cross-section would vary when dragged along $ \nn $, thus defining an ``evolution" along $ \eh $ which would break the invariance of $ \eh $ under the action of the stationary group. A rigorous proof is in Prop. 9.3.1 of \cite{hawking-hellis}. Moreover, if the cross-section $ \crs $ of $ \eh $ are closed manifolds such that the horizon topology is $ \eh \cong \R \times \crs $, then $ \eh $ is a non-expanding horizon (\dref{def:non-expanding-horizon}).

\subsubsection{Topology theorems}
% TODO: add proofs from gourgoulhon

In $ d = 4 $, the topology of the horizon cross-sections can be strongly constrained.

\begin{theorem}{Hawking's topology theorem \cite{hawking-1}}{topology-1}
  Let $ (\man , g) $ be a $ 4 $-dimensional stationary spacetime containing a black hole with event horizon $ \eh $. Let $ \crs $ be a connected component of a complete cross-section of $ \eh $, and assume that:
  \begin{enumerate}
    \item \label{item:htt-1} $ \crs $ is closed (compact withough boundary) and orientable;
    \item \label{item:htt-2} the null dominance condition (\eref{eq:null-dominance-condition}) is satisfied on $ \eh $ for some $ f \in \cfk(\man , \R^+_0) $;
    \item \label{item:htt-3} given an ingoing null normal $ \kk $ to $ \crs $, when displaced into the black hole exterior along $ - \kk $, $ \crs $ becomes a surface with $ \theta_{(\ll)} > 0 $.
  \end{enumerate}
  Then, the cross-section has generically the topology $ \crs \cong \Sn{2} $, or, under specific hypotheses (namely the induced metric $ g \vert_\crs $ being flat), $ \crs \cong \Tn{2} = \Sn{1} \times \Sn{1} $.
\end{theorem}

For General Relativity, \itemref{item:htt-2} with $ f = \Lambda $ is equivalent to assuming $ \Lambda \geq 0 $ and the dominant null energy condition.

This topology theorem is not specific of stationary black hole event horizons, since it relies only on quasi-local properties: in particular, it only requires an hypersurface (not even a null one) sliced by spacelike compact surfaces with $ \theta_{(\nn)} = 0 $ and $ \ld_{\kk} \theta_{(\nn)} \leq 0 $, with a conventional definiton of ``interior" region as that pointed towards by $ \kk $.

Another version of the topology theorem relies on the weaker null convergence condition with global hypotheses, and it excludes the torus topology. In particular, it relies on the notion of global hyperbolicity: given a spacetime $ (\man , g) $, a subset $ U \ssq \man $ is \bctxt{globally hyperbolic} if it admits a \emph{Cauchy hypersurface} (\dref{def:cauchy-hypersurface}).

\begin{theorem}{Chruściel--Wald topology theorem \cite{chrusciel}}{topology-2}
  Let $ (\man , g) $ be a $ 4 $-dimensional asymptotically-flat stationary spacetime containing a black hole with event horizon $ \eh $. Assume that:
  \begin{enumerate}
    \item \label{item:cwtt-1} the null convergence condition (\pref{prop:null-convergence-condition}) is satisfied;
    \item \label{item:cwtt-2} the domain of outer communications $ \doc{\man} $ is globally hyperbolic;
    \item \label{item:cwtt-3} $ \doc{\man} $ contains an achronal asymptotically-flat hypersurface $ \Sigma $ such that $ \eh \cap \Sigma = \crs $, where $ \crs $ is a compact cross-section;
    \item \label{item:cwtt-4} other technical conditions on the spacetime (see \cite{chrusciel});
  \end{enumerate}
  Then, $ \doc{\man} $ is simply connected and any connected component of $ \crs $ is homeomorphic to $ \Sn{2} $.
\end{theorem}

Both \tref{th:topology-1} and \tref{th:topology-2} regard only a connected component of the complete cross-section of the horizon: indeed, there exist stationary (actually static) spacetimes containing a black hole whose event horizon has disconnected complete cross-sections, known as \emph{Majumdar--Papapetrou spacetimes}.
A generalization of \tref{th:topology-1} to spacetimes with $ d > 4 $ is possible, but the result is much weaker and puts less stringent constraints on the topology of cross-sections: this is due to the fact that, for $ d > 4 $, $ \dim_\R \crs > 2 $, but the proof of \tref{th:topology-1} relies on the Gauss--Bonnet theorem\footnotemark which is only valid for $ 2 $-surfaces.

\footnotetext{For a smooth $ 2 $-manifold $ \crs $ with Riemannian metric $ q \in \tns{0}{2}(\crs) $, the Ricci scalar curvature $ \rr_\crs $ associated to $ q $ and the Euler charcteristic $ \chi(\crs) $ are related by:
\begin{equation}
  \int_{\crs} \volf_q \, \rr_\crs = 4 \pi \chi(\crs)
  \label{eq:gauss-bonnet}
\end{equation}}

\subsection{Mass and angular momentum}

For asimptotically-flat stationary spacetimes, it is possible to make the concept of mass well-defined. With the additional requirement of axisymmetry, the concept of angular momentum can be introduced too.

\subsubsection{Weakly-relativistic stationary spacetimes}

\begin{definition}{Weakly-relativistic spacetime}{}
  Let $ (\man , g) $ be a $ d $-spacetime. Then, it is called \bcdef{weakly relativistic} if $ \man $ is diffeomorphic to $ \R^d $ and the metric tensor can be decomposed as:
  \begin{equation}
    g = f + h
    \label{eq:weak-rel-dec}
  \end{equation}
  where $ f $ is a flat Lorentzian metric on $ \man $ and $ h $ is small in the following sense: in any \emph{Minkowskian coordinates}, i.e. $ (x^\alpha) $ on $ \man $ such that $ f_{\alpha \beta} = \eta_{\alpha \beta} $, then $ \abs{h_{\alpha \beta}} \ll 1 $.
\end{definition}

Note that Minkowskian coordinates are not unique, since any Poincaré transformation $ (\Lambda , \aab) \in \ISO{1,d-1} : x^\alpha \mapsto \tensor{\Lambda}{^\alpha_\mu} x^\mu + a^\alpha $ leads to another Minkowskian coordinate frame. Furthermore, the flat metric itself is not unique. Indeed, consider the change of coordinates $ x^\alpha \mapsto y^\alpha = x^\alpha + \zeta^\alpha(x) $, where $ \zeta^\alpha $ are infinitesimal functions, so that:
\begin{equation*}
  \begin{split}
    g'_{\alpha \beta} = \frac{\pa x^\mu}{\pa y^\alpha} \frac{\pa x^\nu}{\pa y^\beta} g_{\mu \nu} = (\delta^\mu_\alpha - \pa_\alpha \zeta^\mu + \smo(\zeta^2)) (\delta^\nu_\beta - \pa_\beta \zeta^\nu + \smo(\zeta^2)) (\eta_{\mu \nu} + h_{\mu \nu})
  \end{split}
\end{equation*}
At first order in $ \zeta $, defining $ \zeta_\alpha \equiv \eta_{\alpha \mu} \zeta^\mu $, then:
\begin{equation*}
  g'_{\alpha \beta} = \eta_{\alpha \beta} + h_{\alpha \beta} - \pa_\alpha \zeta_\beta - \pa_\beta \zeta_\alpha
\end{equation*}
Recasting this expression as $ g' = f' + h' $, where $ f' $ is the metric with components $ f'_{\alpha \beta} = \eta_{\alpha \beta} $ in the $ (y^\alpha) $ coordinate frame, the transformation preserves the flat metric while transforming the ``small perturbation" as:
\begin{equation}
  h'_{\alpha \beta} = h_{\alpha \beta} - \pa_\alpha \zeta_\beta - \pa_\beta \zeta_\alpha
  \label{eq:gauge-freedom}
\end{equation}
This can be regarded as a \emph{gauge freedom} of $ h $, which reflects the freedom in the choice of the flat background metric $ f $ in the decomposition \eref{eq:weak-rel-dec}.

\begin{lemma}{Lorentz gauge}{}
  Let $ (\man , g) $ be a weakly-relativistic spacetime with metric $ g = f + h $, where $ f $ is a flat Lorentzian metric. Then, it is always possible to choose Minkowskian coordinates $ (x^\alpha) $ such that:
  \begin{equation}
    \pa_\mu (\eta^{\mu \nu} \bar{h}_{\alpha \nu}) = 0
    \qquad \qquad
    \bar{h} = h - \frac{1}{2} f \tr h
  \end{equation}
  where $ \tr h \equiv \eta^{\mu \nu} h_{\mu \nu} $.
\end{lemma}

\begin{proofbox}
  \begin{proof}
    By \eref{eq:gauge-freedom} $ \tr h' \equiv \eta^{\mu \nu} h'_{\mu \nu} = \tr h - 2 \pa_\mu \zeta^\mu $, so that $ \bar{h}'_{\alpha \beta} = \bar{h}_{\alpha \beta} - \pa_\alpha \zeta_\beta - \pa_\beta \zeta_\alpha + \eta_{\alpha \beta} \pa_\mu \zeta^\mu $. Thus:
    \begin{equation*}
      \pa_\mu (\eta^{\mu \nu} \bar{h}'_{\alpha \nu}) = \pa_\mu (\eta^{\mu \nu} \bar{h}_{\alpha \nu}) - \dal_f \zeta_\alpha
      \qquad \qquad
      \dal_f \equiv \eta^{\mu \nu} \pa_\mu \pa_\nu
    \end{equation*}
    Therefore, if the Lorentz gauge condition is not fulfilled, it can be enforced with a gauge transformation satisfying $ \dal_f \zeta_\alpha = \pa_\mu (\eta^{\mu \nu} \bar{h}_{\alpha \nu}) $, which always has a solution.
  \end{proof}
\end{proofbox}

Recalling \eref{eq:nlo-einstein} (in this case with flat background), the Einstein tensor reduces to $ \einst = - \frac{1}{2} \dal_f \bar{h} $.

Consider now a weakly-relativistic stationary spacetime $ (\man , g) $ with stationary Killing vector $ \xib $. The Minkowskian coordinates can be chosen so that $ x^0 $ is ignorable, i.e. $ \xib = \pa_0 $. Let $ \Sigma $ be a $ x^0 = \text{constl.} $ hypersurface, and let $ \gamma = f \vert_\Sigma $: the coordinates $ (x^i) $ form a Cartesian coordinate system on the submanifold $ (\Sigma , \gamma) $, since $ \gamma = \delta_{ij} \dd x^i \otimes \dd x^j $, and the d'Alembertian operator $ \dal_f $ reduces to the Laplace--Beltrami operator $ \lap_\gamma $. The Einstein field equations (at linear order in $ h $) then take the form on $ \Sigma $ of $ \frac{1}{2} d (d+1) $ independent Poisson equations:
\begin{equation}
  \lap \bar{h}_{\alpha \beta} = - 16 \pi \gd T_{\alpha \beta}
\end{equation}
The solutions are obtained using the Green's function of the $ (d-1) $-dimensional Laplace--Beltrami operator:
\begin{equation}
  \bar{h}_{\alpha \beta}(x) = \frac{16 \pi \gd}{(d-3) \Omega_{d-2}} \int_\Sigma \dd^{d-1} y \, T_{\alpha \beta}(y) \norm{x - y}^{3-d}
\end{equation}
where the norm is the standard Euclidean norm, since $ \gamma $ is a Euclidean metric. Recall that the surface hyper-area $ \Omega_d $ of $ \Sn{d} $ is given by:
\begin{equation}
  \Omega_d = \frac{2 \pi^{\frac{d+1}{2}}}{\Gamma(\frac{d+1}{2})}
  \qquad \qquad
  \Gamma(z) \defeq \int_0^{\infty} \dd t \, e^{-t} t^{z-1} \ , \ z \in \C : \Re z > 0
\end{equation}

\begin{proposition}{Asymptotic metric of weakly-relativistic system}{}
  Let $ (\man , g) $ be a weakly-relativistic stationary $ d $-spacetime, with $ d \geq 4 $, obeying the Einstein field equations \eref{eq:ef-eq} with $ \Lambda = 0 $ and an EMT $ T $ with compact support (the \emph{source}), where the energy density dominates over spatial stresses (weakly-relativistic matter). In the Lorentz gauge, there exists a coordinate frame $ (x^\alpha) $ such that the asymptotic behaviour of the metric tensor for $ r^2 \equiv \delta_{ij} x^i x^j \ra \infty $ is:
  \begin{equation}
    g_{00} = -1 + \alpha_d \frac{M}{r^{d-3}} + \smo(r^{1-d})
    \label{eq:asymptotic-00}
  \end{equation}
  \begin{equation}
    g_{0i} = \frac{d-2}{2} \alpha_d \frac{J_{ij} x^j}{r^{d-1}} + \smo(r^{1-d})
    \label{eq:asymptotic-0i}
  \end{equation}
  \begin{equation}
    g_{ij} = \left( 1 + \frac{\alpha_d}{d-3} \frac{M}{r^{d-3}} \right) \delta_{ij} + \smo(r^{2-d})
  \end{equation}
  where $ \alpha_d \equiv 16 \pi / ((d-2) \Omega_{d-2}) $. The \bcprop{mass} and the \bcprop{angular momentum} of the source are:
  \begin{equation}
    M \defeq \int_\Sigma \dd^{d-1} x \, T_{00}(x)
    \qquad \qquad
    J_{ij} \defeq \int_\Sigma \dd^{d-1} x \left( x^j T_{0i}(x) - x^i T_{0j}(x) \right)
    \label{eq:wrs-mass}
  \end{equation}
  where $ \Sigma $ is a $ x^0 = \text{const.} $ hypersurface ($ M $ and $ J_{ij} $ are independent of $ \Sigma $, i.e. of $ x^0 $). Such a coordinate frame is called the \bcprop{central source rest-frame}, since their origin is in the center of mass of the source, in the sense that:
  \begin{equation}
    \int_\Sigma \dd^{d-1} x \, T_{0i}(x) = 0
    \qquad \qquad
    \int_\Sigma \dd^{d-1} x \, x^i T_{00}(x) = 0
  \end{equation}
  The weakly-relavistic nature of the source is expressed in the $ (x^\alpha) $ frame as $ T_{00} \gg \abs{T_{ij}} $.
\end{proposition}
% TODO: add proof from gourgoulhon

\subsubsection{Komar mass}

To generalize the concept of mass to strongly-relativistic spacetimes in a coordinate-independent way, consider the fact that, in Newtonian gravity, the mass is provided by Gauss's law for the gravitational field $ \ve{g} = - \grad \Phi $:
\begin{equation*}
  M = \frac{1}{4\pi} \int_\surf \dd \ve{S} \cdot \grad \Phi
\end{equation*}
where $ \surf $ is any closed surface surrounding the body. The arbitrariness of $ \surf $ is possible thanks to the Gauss--Ostrogradsky theorem and to the Laplace equation $ \lap \Phi = 0 $ (outside the source). Analogously, note that the mass $ M $ appears as the coefficient of the dominant $ r^{3-d} $ term in the asymptotic expansion of $ g_{00} $ in \eref{eq:asymptotic-00}, so it could be recovered considering the flux of $ \pa_r g_{00} $ through a $ (d-2) $-sphere at large $ r $.

In a coordinate-invariant way, it is more appropriate to consider the stationary Killing vector field $ \xib \in \xm $, and in particular its metric-dual $ 1 $-form $ \xi^\flat \in \df{1} $, since in adapted coordinates $ \xi^\alpha = (1 , 0 , \dots , 0) $ and $ \xi_\alpha = g_{0\alpha} $. Then, the concept of flux is naturally generalized by Komar charges (\secref{sssec:komar}), so that, assuming $ \xib $ to be normalized as $ \braket{\xib , \xib}_g \ra -1 $ on the asymptotic boundary, the \bctxt{Komar mass} over a spacelike closed $ (d-2) $-surface $ \surf \ssq \man $ is defined as:
\begin{equation}
  M_\surf \defeq \frac{d-2}{d-3} Q_\text{K}[\xib ; \surf] = - \frac{d-2}{d-3} \frac{1}{16\pi \gd} \int_\surf \star \dd \xi^\flat
  \label{eq:komar-mass}
\end{equation}
The orientation of $ \surf $ is set by the $ (d-2) $-form $ \volf_\surf \equiv \ip_{\ssb} \ip_{\nb} \volf_g $, where $ \nb $ is any future-directed unit timelike vector field normal to $ \surf $ and $ \ssb $ is a unit spacelike vector field normal to $ \nb $, orthogonal to $ \nb $ and directed towards the ``exterior" of $ \surf $ in the following sense: assuming that $ \surf $ can be embedded into a spacelike hypersurface $ \Sigma \ssq \man $ that extends towards the asymptotic flat end of $ (\man , g) $, then for any point $ p \in \surf $ there exists a curve entirely lying in $ \Sigma $, starting from $ p $ with tangent vector $ \ssb_p $ and extending to the asymptotic flat end withour crossing $ \surf $ again.

Instead of the integral of a $ (d-2) $-form over a $ (d-2) $-surface, it is possible to express the Komar mass as a flux integral, i.e. the integral of a $ 2 $-form contracted with an ``area element" normal to the $ (d-2) $-surface. In particular, given a spacelike $ (d-2) $-surface $ \surf \ssq \man $, the \bctxt{area element normal bivector} to $ \surf $ is the antisymmetric $ (2,0) $-tensor field defined on $ \surf $ as:
\begin{equation}
  \dd \ve{S} \defeq (\ssb \wedge \nb) \sqrt{\mt{q}} \, \dd^{d-2} x
  \label{eq:area-element}
\end{equation}
where $ \nb $ and $ \ssb $ are the unit timelike and spacelike vector fields normal to $ \surf $ defined above (in particular, at each $ p \in \surf $, an orthonormal basis of the timelike plane $ T_p^\perp \surf $ normal to $ \surf $ is $ (\nb_p , \ssb_p) $), $ \mt{q} \equiv \det q_{ab} $ is the determinant of the induced (Riemannian) metric $ q = g \vert_\surf $, and $ (x^a) = (x^1 , \dots , x^{d-2}) $ is a coordinate frame on $ \surf $ such that $ (\pa_1 , \dots , \pa_{d-2}) $ is a right-handed basis for the orientation $ \volf_\surf $. The \emph{exterior product} of vectors is analogous to that of forms: $ \ssb \wedge \nb \defeq \ssb \otimes \nb - \nb \otimes \ssb $. The bivector $ \dd \ve{S} $ can be shown to be independent of the choice of coordinates $ (x^a) $ on $ \surf $ and of orthonormal basis $ (\nb , \ssb) $ of $ T_p^\perp \surf $.
% TODO: add proof from Prop. 5.9 of gourgoulhon

\begin{lemma}{Flux integral}{flux}
  Let $ (\man , g) $ be a $ d $-spacetime, let $ \omega \in \df{2} $ and let $ \surf \ssq \man $ be a spacelike $ (d-2) $-surface. Then:
  \begin{equation}
    \int_\surf \star \omega = \frac{1}{2} \int_\surf \dd S^{\mu \nu} \omega_{\mu \nu}
    \label{eq:flux}
  \end{equation}
  where $ \dd \ve{S} $ is the area element normal bivector to $ \surf $ defined by \eref{eq:area-element}.
\end{lemma}
% TODO: add proof from Lemma 5.10 of gourgoulhon

Using \eref{eq:flux}, since $ (\dd \xi^\flat)_{\mu \nu} = 2 \na_\mu \xi_\nu $ (by metric compatibility and Killing's equation), the Komar mass can be rewritten as a flux integral:
\begin{equation}
  M_\surf = - \frac{d-2}{d-3} \frac{1}{16 \pi \gd} \int_\surf \dd S^{\mu \nu} \na_\mu \xi_\nu
  \label{eq:komar-mass-flux}
\end{equation}
Alternatively, without using covariant derivatives:
\begin{equation}
  M_\surf = - \frac{d-2}{d-3} \frac{1}{16 \pi \gd} \int_\surf \dd^{d-2} x \, \sqrt{\mt{q}} \, \pa_\mu \xi_\nu (s^\mu n^\nu - n^\mu s^\nu)
  \label{eq:komar-mass-flux-alt}
\end{equation}

\begin{example}{Komar mass of Schwarzschil black hole}{}
  Consider the Schwarzschild spacetime $ (\man , g) $ in ingoing Eddington--Finkelstein coordinates $ (t,r,\vartheta,\varphi) $, so that the metric is \eref{eq:in-edd-fink}. The stationary Killing vector is $ \xib = \pa_t $, and choose $ \surf $ to be a $ 2 $-sphere defined by $ (t,r) = \text{const.} $, so that the coordinates spanning $ \surf $ are $ (x^a) = (\vartheta , \varphi) $. For the orthonormal basis of $ T_p^\perp \surf $, choose:
  \begin{equation*}
    \nb = N^{-1} \frac{\pa}{\pa t} - \frac{2mN}{r} \frac{\pa}{\pa r}
    \qquad
    \ssb = N \frac{\pa}{\pa r}
    \qquad \qquad
    N \equiv \left( 1 + \frac{2m}{r} \right)^{-1/2}
  \end{equation*}
  which are easily checked to satisfy $ \braket{\nb , \nb}_g = -1 $, $ \braket{\ssb , \ssb}_g = 1 $ and $ \braket{\nb , \ssb}_g = 0 $. Moreover, $ \braket{\nb , \pa_\vartheta}_g = \braket{\nb , \pa_\varphi}_g = 0 $ and $ \braket{\ssb , \pa_\vartheta}_g = \braket{\ssb , \pa_\varphi}_g = 0 $, which proves that $ \nb $ and $ \ssb $ are normal to $ \surf $, and $ \ssb $ points to the exterior of $ \surf $ (since $ N > 0 $), so that $ (\nb , \ssb , \pa_\vartheta , \pa_\varphi) $ is a right-handed basis. To evaluate the Komar mass, compute:
  \begin{equation*}
    \xi_\alpha = g_{\alpha t} = \left( -1 + \frac{2m}{r} , \frac{2m}{r} , 0 , 0 \right)
  \end{equation*}
  which means that the only non-zero derivatives $ \pa_\mu \xi_\nu $ are $ \pa_r \xi_t = \pa_r \xi_r = - 2m / r^2 $. The integrand in \eref{eq:komar-mass-flux-alt} reduces to:
  \begin{equation*}
    \pa_\mu \xi_\nu (s^\mu n^\nu - n^\mu s^\nu) = \pa_r \xi_t (s^r n^t - n^r s^t) = - \frac{2m}{r^2} (N \cdot N^{-1} + 2mN \cdot 0) = - \frac{2m}{r^2}
  \end{equation*}
  From \eref{eq:in-edd-fink}, the induced metric on $ \surf $ is $ \dd s^2 = r^2 \dd \Omega_2^2 $, so that $ \sqrt{\mt{q}} = r^2 \sin \vartheta $ and:
  \begin{equation*}
    M_\surf = - \frac{1}{8\pi \gn} \int_\surf \dd \vartheta \dd \varphi \, r^2 \sin \vartheta \left( - \frac{2m}{r^2} \right) = \frac{m}{4\pi \gn} \int_\surf \dd \vartheta \dd \varphi = \frac{m}{\gn}
  \end{equation*}
  Note that $ M_\surf $ does no depend on $ r $, i.e. it does not depend on the choice of the sphere $ \surf $.
\end{example}

In general, it can be shown that, for a stationary asymptotically-flat spacetime, the mass in \eref{eq:asymptotic-00} is precisely the Komar mass over a sphere at constant $ (x^0 , r) $, with $ r \ra \infty $: in particular, if the spacetime is weakly relativistic, the Komar mass is given by \eref{eq:wrs-mass} (integral over $ \surf $).
% TODO: add proof from Prop. 5.11 of gourgoulhon
Another useful expression for the Komar mass can be obtained in the form of a volume integral.

\begin{lemma}{Flux and volume integrals}{volume-integral}
  Let $ (\man , g) $ be an asymptotically-flat $ d $-spacetime, and let $ \Sigma $ be a compact spacelike hypersurface bounded by two closed $ (d-2) $-surfaces, an ``internal" one $ \surf_\text{int} $ and an ``external" one $ \surf_\text{ext} $, where the exterior direction is that of the asymptotic flat end of $ (\man , g) $. Let $ \nb $ be the future-directed unit normal to $ \Sigma $, and define the \bclemma{volume element normal vector} to $ \Sigma $ as:
  \begin{equation}
    \dd \ve{V} \defeq - \nb \sqrt{\gamma} \, \dd^{d-1} x \equiv - \nb \volf_\Sigma
  \end{equation}
  where $ (x^i) = (x^1 , \dots , x^{d-1}) $ is a coordinate frame on $ \Sigma $ and $ \gamma \equiv \det \gamma_{ij} $ is the determinant of the induced metric on $ \Sigma $. Then, for any $ \omega \in \df{2} $:
  \begin{equation}
    2 \int_\Sigma \dd V^\mu \na^\nu \omega_{\mu \nu} = \int_{\surf_\text{ext}} \dd S^{\mu \nu} \omega_{\mu \nu} - \int_{\surf_\text{int}} \dd S^{\mu \nu} \omega_{\mu \nu}
  \end{equation}
  where $ \dd S^{\mu \nu} $ is defined by \eref{eq:area-element}, with $ \ssb $ pointing to the asymptotic flat end of $ (\man , g) $ for both $ \surf_\text{ext} $ and $ \surf_\text{int} $.
\end{lemma}
% TODO: add proof from Lemma 5.12 of gourgoulhon

Recalling \eref{eq:killing-riemann}, for Killing vector fields $ \na_\mu \na^\alpha \xi^\mu = \tensor{\rr}{^\alpha_\mu} \xi^\mu $, so, considering the $ 2 $-form $ \dd \xi^\flat $, which has components $ (\dd \xi^\flat)_{\mu \nu} = 2 \na_\mu \xi_\nu $, and applying \lref{lemma:volume-integral}, for an asymptotically-flat spacetime $ (\man , g) $:
\begin{equation}
  M_{\surf_\text{ext}} = M_{\surf_\text{int}} + \frac{d-2}{d-3} \frac{1}{8\pi \gn} \int_\Sigma \volf_\Sigma \, \ricc(\xib , \nb)
  \label{eq:komar-mass-cons}
\end{equation}
where $ \Sigma $ is a compact spacelike hypersurface like in \lref{lemma:volume-integral}. The same holds if $ \Sigma $ is bounded by a single closed $ (d-2) $-surface, provided that the Komar mass of the other (non-existing) surface is set to zero. If $ (\man , g) $ is ruled by General Relativity, then in vacuum $ \ricc \equiv 0 $, which means that the Komar mass $ M_\surf $ is independent of the $ (d-2) $-surface $ \surf $ chosen. In particular, assuming that the EMT of the matter distribution vanishes for $ r \geq r_0 $, for some $ r_0 \in \R^+ $, then the Komar mass of any $ (d-2) $-surface $ \surf_r $ defined by $ (x^0 , r) = \text{const.} $ with $ r \geq r_0 $ is equal to that computed on a $ (d-2) $-sphere at infinity: $ M_\infty \equiv M_{\Sn{d-2}_\infty} $.

The Komar mass can be regarded as the \emph{total mass} of a spacetime, with contributions from both matter fields and the gravitational field.

% TODO: add observation on the equivalence between Komar mass and ADM mass

\subsubsection{Komar angular momentum}

As a stationary Killing vector determines a Komar mass, an axisymmetric Killing vector determines a Komar angular momentum.

\begin{definition}{Axisymmetric spacetime}{}
  A $ d $-spacetime $ (\man , g) $ is called \bcdef{axisymmetric} if it is invariant under the action of $ \SOn{2} \cong \Un{1} $ and the orbits of the group action are everywhere spacelike curves.
\end{definition}

Since $ \dim_\R \SOn{2} = 1 $, it is generated by a single Killing vector field $ \pb \in \xm $, which is unique up to a rescaling constant: since the orbits of $ \SOn{2} $ are closed curves, the rescaling constant is fixed requiring that the group parameter $ \varphi $ associated with $ \pb $ has the standard periodicity $ 2\pi $. If the set of points of $ \man $ left invariant under $ \SOn{2} $ is non-empty, it is called the \bctxt{rotation axis}: equivalently, it is the set of points $ p \in \man : \pb_p = \bs{0} $. Usually, the rotation axis is a $ (d-2) $-dimensional submanifold of $ \man $, and it is timelike in the regione where $ \SOn{2} $ orbits are spacelike.

Given an axisymmetric $ d $-spacetime $ (\man , g) $ with $ d \geq 3 $ and axisymmetric Killing vector field $ \pb \in \xm $, the \bctxt{Komar angular momentum} over a spacelike closed $ (d-2) $-surface $ \surf \ssq \man $ is defined as:
\begin{equation}
  J_\surf \defeq - Q_\text{K}[\pb , \surf] = \frac{1}{16 \pi \gd} \int_\surf \star \dd \psi^\flat
  \label{eq:komar-angular-momentum}
\end{equation}
Notice the absence of the dimensional factor, when compared to \eref{eq:komar-mass}: this independence of \eref{eq:komar-angular-momentum} from the spacetime dimension is due to the $ \SOn{2} $ action effective only in a $ 2 $-dimensional space, which is the orthogonal complement of the rotation axis. Moreover, $ J_\surf $ is defined for $ d \geq 3 $, while $ M_\text{surf} $ only for $ d \geq 4 $.

Since \eref{eq:komar-mass-flux} and \eref{eq:komar-mass-flux-alt} only rely on $ \xib $ being a Killing vector, they remain valid for $ \pb $ (with the appropriate prefactor):
\begin{equation}
  J_\surf = \frac{1}{16 \pi \gd} \int_\surf \dd S^{\mu \nu} \na_\mu \psi_\nu
  \label{eq:komar-angular-momentum-flux}
\end{equation}
\begin{equation}
  J_\surf = \frac{1}{16 \pi \gd} \int_\surf \dd^{d-2} x \, \sqrt{\mt{q}} \, \pa_\mu \psi_\nu (s^\mu n^\nu - n^\mu s^\nu)
  \label{eq:komar-angular-momentum-flux-alt}
\end{equation}
Moreover, it can be shown that, for a stationary axisymmetric asymptotically-flat spacetime, if asymptotically Minkowskian coordinates adapted to $ \pb $ are used, so that $ \pb = - x^2 \pa_{x^1} + x^1 \pa_{x^2} $ (i.e. the rotation axis is the $ (d-2) $-surface $ (x^1 , x^2) = (0,0) $), then the coefficient $ J_{12} $ in \eref{eq:asymptotic-0i} is precisely the Komar angular momentum over a sphere at constant $ (x^0 , r) $, with $ r \ra \infty $: in particular, if the spacetime is weakly relativistic, the Komar angular momentum is given by \eref{eq:wrs-mass} ($ J_{12} $ over $ \surf $).
% TODO: add proof from Prop. 5.19 of gourgoulhon
\eref{eq:komar-mass-cons} too can be imported to the Komar angular momentum:
\begin{equation}
  J_{\surf_\text{ext}} = J_{\surf_\text{int}} + \frac{1}{16 \pi \gd} \int_\Sigma \dd^{d-1} x \, \sqrt{\gamma} \left( -2 \tensor{\rr}{^\alpha_\mu} \psi^\mu \right) n_\alpha
\end{equation}
where $ \Sigma $ is a compact spacelike hypersurface bounded by two closed $ (d-2) $-surfaces, $ \nb $ is the future directed unit normal to $ \Sigma $ and $ \gamma \equiv \det \gamma_{ij} $ is the determinant of the induced metric on $ \Sigma $. Therefore, in a vacuum spacetime, the Komar angular momentum too is independent of the $ (d-2) $-surface on which it is computed.

The Komar angular momentum can be regarded as the \emph{total angular momentum} of a spacetime, with contributions from both matter fields and the gravitational field.















































\newpage
\section{Schwarzschild black holes}
\label{sec:bh-schwarzschild}

% TODO: must rewiew

The simplest black hole solution is the Schwarzschild metric:
\begin{equation}
  \dd s^2 = - \left( 1 - \frac{2GM}{r} \right) \dd t^2 + \left( 1 - \frac{2GM}{r} \right)^{-1} \dd r^2 + r^2 \left( \dd \vartheta^2 + \sin^2 \vartheta\, \dd \varphi^2 \right)
  \label{eq:schw-metric}
\end{equation}
This is a special case of the general metric in \secref{ssec:dS-space} with $ f(r)^2 = 1 - 2GM / r $: it solves the vacuum Einstein equations $ R_{\mu \nu} = 0 $.

This solution depends on a single parameter $ M $, which is identified as the mass of the black hole, as highlighted in \secref{sec:schwarz-sol}. The same can be shown via Komar integrals: note that the Schwarzschild metric admits a timelike Killing vector field $ K = \pa_t $, with dual 1-form $ K^\flat = g_{00} \dd t $ and Komar charge:
\begin{equation*}
  F = \dd K^\flat = - \frac{2GM}{r^2} \dd r \wedge \dd t
  \quad \implies \quad
  M_{\text{Komar}} = - \frac{1}{8\pi G} \int_{\Sn{2}_R} \star\, F = M
\end{equation*}
for all $ R > 2GM $ (radius of the horizon). Although the 2-form $ F $ obeys the vacuum Maxwell equations $ \dd \star F = 0 $, thus one would expect the associated Komar charge to vanish as there is no current, $ M_{\text{Komar}} = M \neq 0 $: this is possible because the mass of the black hole is localized at the origin, where $ F $ diverges. Moreover, the Schwarzschild solution is mathematically valid for $ M \in \R $: the solution $ M = 0 $ is just Minkowski spacetime, while $ M < 0 $ will be shown to be unphysical.

\subsection{Horizon}
\label{sec:schwarz-horizon}

The Schwarzschild metric diverges at $ r = 0 $ and $ r = 2GM \equiv \sr $. To distinguish between true singularities and coordinate singularities, the usual way is defining a coordinate-independent scalar quantity and studying it at divergence points.

Due to the vacuum field equations, the simplest scalars $ R $ and $ R_{\mu \nu} R^{\mu \nu} $ both vanish: the next simplest curvature scalar is the \bctxt{Krentschmann scalar} $ R^{\mu \nu \rho \sigma} R_{\mu \nu \rho \sigma} $, which for the Schwarzschild metric is:
\begin{equation}
  R^{\mu \nu \rho \sigma} R_{\mu \nu \rho \sigma} = \frac{48 G^2 M^2}{r^6}
\end{equation}
At $ r = 2GM $, the Krentschmann scalar is $ \sim 1 / (GM)^4 $, suggesting that this is only a coordinate singularity: it turns out that is the case, and indeed the surface $ r = \sr $ is known as the \bctxt{event horizon} of the black hole. Note that, interestingly, heavier black holes have smaller curvature at the horizon.

On the contrary, $ r = 0 $ is a true singularity, simply known as the \bctxt{singularity} of the (classical) black hole.

\subsubsection{Near horizon limit}

To study the spacetime in the vicinity of the horizon, set $ r = 2GM + \eta $, with $ \abs{\eta} \ll 2GM $. Moreover, consider $ \eta \in \R^+ $, thus restricting to the spacetime just outside the horizon. The metric then becomes, at first order in $ \eta $:
\begin{equation*}
  \dd s^2 = - \frac{\eta}{2GM} \dd t^2 + \frac{2GM}{\eta} \dd \eta^2 + (2GM)^2 \dd \Omega_2^2
\end{equation*}
Spacetime is thus decomposed in a direct product of $ \Sn{2}_{2GM} $ and a 2-dimensional Lorentzian manifold. Focusing on the latter, set $ \rho^2 = 8GM \eta $, so that:
\begin{equation}
  \dd s^2 = - \left( \frac{\rho}{4GM} \right)^2 \dd t^2 + \dd \rho^2
\end{equation}
This is called \bctxt{Rindler metric} and it is, in fact, just Minkowski spacetime:
\begin{equation*}
  T = \rho \sinh \left( \frac{t}{4GM} \right)
  \qquad
  X = \rho \cosh \left( \frac{t}{4GM} \right)
  \quad \implies \quad
  \dd s^2 = - \dd T^2 + \dd X^2
\end{equation*}
These are precisely the coordinates \eref{eq:km-coord} experienced by an observer undergoing constant acceleration $ a = 1 / (4GM) $: this makes sense, as an observer sitting at constant $ \rho $, i.e. constant $ r $, must accelerate in order not to fall into the black hole.

The outside of the black hole is $ (\rho,t) \in \R^+ \times \R $, corresponding to $ X > \abs{T} $. The horizon $ \rho = 0 $, instead, is mapped to the origin $ T = X = 0 $ of Minkowski space. However, note that the time coordinate is undefined at the origin since $ g_{00} = 0 $: scaling $ t \rightarrow \infty $ and $ \rho \rightarrow 0 $ keeping $ \rho e^{\pm t / 4GM} $ fixed makes it clear that the horizon actually corresponds to the lines:
\begin{equation*}
  r = 2GM \quad \implies \quad X = \pm T
\end{equation*}
Therefore, the event horizon is not a timelike surface (like the surface of a star), but a null surface.

Although the outside of the event horizon is only described by $ X > \abs{T} $, it is nonetheless possible to extend the metric to the whole Minkowski space $ X,T \in \R $: this makes it clear that the event horizon is not a real singularity, as nothing particular happens at $ X = \pm T $ (see \figref{fig:rindler}). Zooming out, however, the peculiar properties of the horizon start to emerge from a global perspective.

\begin{figure}[b]
  \centering
  \includegraphics[width = 0.35 \textwidth]{rindler.png}
  \caption{Rindler spacetime, with null lines in black and a line at constant $ r > 2GM $ in red.}
  \label{fig:rindler}
\end{figure}

\subsection{Eddington--Finkelstein coordinates}
\label{ssec:eddington-finkelstein-coordinates}

Both the Schwarzschild coordinates and the Rindler coordinates do not cover the whole manifold. To find a set of coordinates that does, first introduce a new radial coordinate:
\begin{equation}
  \dd r_*^2 = \left( 1 - \frac{2GM}{r} \right)^{-2} \dd r^2
  \quad \implies \quad
  r_* = r + 2GM \log \left( \frac{r - 2GM}{2GM} \right)
\end{equation}
This is the \bctxt{Regge--Wheeler radial coordinate} and it maps the region outside the horizon $ r \in (2GM, \infty) $ to $ r_* \in \R $: as an observer approaches the horizon, $ r $ changes increasingly slow varying $ r_* $, since $ \frac{\dd r}{\dd r_*} \rightarrow 0 $ as $ r \rightarrow 2GM $. This coordinate is suited to describe the path of lightrays travelling in the radial direction:
\begin{equation*}
  \dd s^2 = 0
  \quad \implies \quad
  \frac{\dd r}{\dd t} = \pm \left( 1 - \frac{2GM}{r} \right)
  \quad \implies \quad
  \frac{\dd r_*}{\dd t} = \pm 1
  \quad \implies \quad
  t \pm r_* = \text{const.}
\end{equation*}
These null radial geodesics correspond to an ingoing lightray for the positive sign and an outgoing lightray for the negative sign ($ r_* $ must decrease/increase as $ t $ increases). Next, introduce a pair of null coordinates:
\begin{equation}
  v = t + r_*
  \qquad
  u = t - r_*
\end{equation}
These allow to extend the Schwarzschild solution beyond the horizon.

\subsubsection{Ingoing coordinates}

Replace $ t $ with $ t = v - r_*(r) $:
\begin{equation*}
  \dd t = \dd v - \left( 1 - \frac{2GM}{r} \right)^{-1} \dd r
\end{equation*}
The Schwarzschild metric then becomes, in coordinates $ (v,r) $:
\begin{equation}
  \dd s^2 = - \left( 1 - \frac{2GM}{r} \right) \dd v^2 + 2 \dd v\,\dd r + r^2 \dd \Omega_2^2
  \label{eq:ingoing-coord}
\end{equation}
This is the Schwarzschild black hole in \bctxt{ingoing Eddington--Finkelstein coordinates}. Note that there is no more singularity at $ r = 2GM $; however, the $ \dd v^2 $ term vanishes at the horizon and becomes negative for $ r < 2GM $. Nonetheless, the metric is still non-degenerate thanks to the cross-term:
\begin{equation*}
  \det g = \det
  \begin{bmatrix}
    - (1 - \frac{2GM}{r}) & 1 & 0 & 0 \\
    1 & 0 & 0 & 0 \\
    0 & 0 & r^2 & 0 \\
    0 & 0 & 0 & r^2 \sin^2 \vartheta
  \end{bmatrix}
  = - r^4 \sin^2 \vartheta
\end{equation*}
Thus, the metric is non-degenerate and with Lorentzian signature for all values of $ r \in \R_+ $ (except for the known poles of $ \Sn{2} $): this means that the radial coordinate can be extended past the horizon. The globally-timelike Killing vector field $ K = \pa_t $ of the Schwarzschild metric is retained in the ingoing Eddington--Finkelstein coordinates as $ K = \pa_v $, however now it is no longer globally-timelike: it remains so outside the horizon, where $ g_{vv} < 0 $, but become spacelike inside it, where $ g_{vv} > 0 $. This means that the full black hole geometry is not time-independent.

Recall that ingoing lightrays follow geodesics given by $ v = \text{const.} $, while outgoing lightrays follow $ u = t - r_* = \text{const.} $. As a function of $ (v,r_*) $, the latter reads $ v = 2r_* + \text{const.} $, so in $ (v,r) $ coordinates:
\begin{equation*}
  v =  2r + 4GM \log \left( \frac{r - 2GM}{2GM} \right) + \text{const.}
\end{equation*}
This clearly has domain $ r \in (2GM, \infty) $, therefore the Regge--Wheeler radial coordinate needs a redefinition:
\begin{equation*}
  r_* = r + 2GM \log \abs{\frac{r - 2GM}{2GM}}
\end{equation*}
This means that now $ r_* $ is a multi-valued function: $ r \in (2GM, \infty) $ is mapped to $ r_* \in \R $, while $ r \in (0, 2GM) $ is mapped to $ r_* \in \R_- $, with the singularity at $ r_* = 0 $. Outgoing lightrays inside the horizon will then follow geodesics given by:
\begin{equation*}
  v = 2r + 4GM \log \left( \frac{2GM - r}{2GM} \right) + \text{const.}
\end{equation*}
Outgoing null geodesics at the horizon are easily found: the $ \dd v^2 $ term in the metric \eref{eq:ingoing-coord} vanishes and the surface $ r = 2GM $ is itself a null geodesic (as event horizons are always null surfaces).

\begin{figure}
  \centering
  \includegraphics[width = 0.50 \textwidth]{ef-in.png}
  \caption{Finkelstein diagram in ingoing coordinates: ingoing null geodesics are shown in red, outgoing ones in blue.}
  \label{fig:ef-in}
\end{figure}

All of this information about lightrays can be pictured in a \bctxt{Finkelstein diagram}, which shows null geodesics in a radial-temporal diagram. Since $ r_* $ is multi-valued, $ r $ is chosen as radial coordinate: as a consequence, one needs to define a new temporal coordinate $ t_* : v = t + r_* = t_* + r $, i.e. $ t_* = v - r $. The Finkelstein diagram in ingoing coordinates is shown in \figref{fig:ef-in}: ingoing null geodesics are $ v = \text{const.} $, i.e. $ t_* = \text{const.} - r $, hence travelling at 45°, while outgoing ones are given by the previously found expressions. Note that only outgoing null geodesics outside of the horizon effectively ``go out", since $ r \rightarrow \infty $ as $ t \rightarrow \infty $ (so $ t_* \rightarrow \infty $): those inside the horizon move inexorably towards the curvature singularity at $ r = 0 $ as time passes, and each one of them reaches it in a finite time $ t_* $. The boundary between these two regions are the null geodesics which run along the horizon at $ r = 2GM $.

This analysis can be extended to massice particles: they must move along timelike geodesics inside their future lightcone, which is determined by a pair of ingoing and outgoing future-pointing null geodesics. These future lightcones are also shown in \figref{fig:ef-in}: the name ``black hole" captures the causal structure of the spacetime, as neither massive particles nor lightrays can escape outside of the event horizon, thus an outside observer can know nothing of the inside of the black hole.

Consider two observers, one at constant $ r > 2GM $ and the other moving towards the singularity: as the second approaches the event horizon, light signals emitted by it will take longer and longer to reach the first, and they will appear progressively more redshifted due to the increasingly strong gravitational well. Thus, the outside observer will see the ingoing observer forever, but progressively more slowed down and redshifted, and will never see it passing the event horizon.

\subsubsection{Outgoing coordinates}

A different extension of the exterior of the Schwarzschild black hole is obtained by replacing $ t $ with $ u = t - r_* $, i.e. $ t = u + r_*(r) $:
\begin{equation}
  \dd s^2 = - \left( 1 - \frac{2GM}{r} \right) \dd u^2 -2 \dd u\,\dd r + r^2 \dd \Omega_2^2
  \label{eq:outgoing-metric}
\end{equation}
This is the Schwarzschild black hole in \bctxt{outgoing Eddington--Finkelstein coordinates}. Once again, the metric is smooth and non-degenerate at the horizon, thus it is possible to extend $ r \in \R_+ $. However, now $ r \in (0,2GM) $ describe a different part of spacetime from the analogous region in ingoing coordinates. To see this, set coordinates $ (r,t_*) : t_* = u + r $ and draw the Finkelstein diagram (\figref{fig:ef-out}): outgoing null geodesics hence travel at 45°, effectively ``going out" regardless of their position with respect to the horizon. On the contrary, ingoing null geodesics now do not ``go in": those starting outside of the horizon pile up at the horizon itself, unable to cross it, while those starting inside the horizon move outwards and pile up at it, still unable to cross it.

\begin{figure}
  \centering
  \includegraphics[width = 0.50 \textwidth]{ef-out.png}
  \caption{Finkelstein diagram in outgoing coordinates: ingoing null geodesics are shown in red, outgoing ones in blue.}
  \label{fig:ef-out}
\end{figure}

Studying lightcones, it is easy to see that massive particles inside the horizon are inexorably expelled outside of it in a finite amount of time $ t_* $. This is clearly a very different physics from a black hole: the metric \eref{eq:outgoing-metric} is that of a \bctxt{white hole}, an object which expells any matter inside. A white hole can be viewed as the time reversal of a black hole: this is to be traced to the negative sign in the cross-term and can be seen by flipping \figref{fig:ef-in} upside-down, getting \figref{fig:ef-in}.

White holes are perfectly acceptable solutions to the Einstein field equations, and their existence is to be expected from the time-reversal invariance of the field equations. However, they're not physically relevant, as they cannot be formed from collapsing matter.

\subsection{Kruskal spacetime}

To better understand how the region parametrized by $ r \in (0, 2GM] $ corresponds to two different parts of spacetime, one needs to introduce coordinates which cover the entire spacetime, including both black and white holes.

First, rewrite the Schwarzschild metric in $ (u,v) $ coordinates:
\begin{equation*}
  \dd s^2 = - \left( 1 - \frac{2GM}{r} \right) \dd u\,\dd v + r^2 \dd \Omega_2^2
\end{equation*}
where $ r = r(u - v) $. There's a degeneracy at $ r = 2GM $ which can be solved introducing the \bctxt{Kruskal--Szekeres coordinates}:
\begin{equation}
  U = - \exp \left( - \frac{u}{4GM} \right)
  \qquad \qquad
  V = \exp \left( \frac{v}{4GM} \right)
\end{equation}
Both $ U $ and $ V $ are null coordinates, and the exterior of the balck hole is parametrized by $ U < 0, V > 0 $. The metric becomes:
\begin{equation}
  \dd s^2 = - \frac{32 (GM)^3}{r} e^{-\frac{r}{2GM}} \dd U\,\dd V + r^2 \dd \Omega_2^2
  \label{eq:kruskal-metric}
\end{equation}
This is known as the metric of \textit{Kruskal spacetime}. The function $ r = r(U,V) $ can be found inverting:
\begin{equation}
  UV = - \exp \left( \frac{r_*}{2GM} \right) = \frac{2GM - r}{2GM} \exp \left( \frac{r}{2GM} \right)
  \label{eq:kruskal-coord-1}
\end{equation}
On the other hand, the function $ t = t(U,V) $ can be found inverting:
\begin{equation}
  \frac{U}{V} = - \exp \left( - \frac{t}{2GM} \right)
  \label{eq:kruskal-coord-2}
\end{equation}
The original Schwarzschild metric, which covers $ U < 0, V > 0 $, can now be extended to $ (U,V) \in \R^2 $, without any divergence at $ r = 2GM $.

\begin{proposition}{}{}
  Kruskal spacetime is the \bcprop{maximal extension} of the Schwarzschild metric.
\end{proposition}

\begin{proofbox}
  \begin{proof}
    To check whether a given spacetime can be further extended, one needs to look at all possible geodesics: if they are defined for infinite affine parameter, then they can escape to infinity; however, if they alt at some finite affine parameter, that is either a coordinate singularity, hinting that the spacetime can in fact be extended, or a true singularity. A spacetime is maximally extended if any geodesics that alts at a finite affine parameter does so at a true singularity, and that is the case for Kruskal spacetime.
  \end{proof}
\end{proofbox}

The extension process allows to solve the field equations in some particular region (open subset) of spacetime: being the metric components real analytic functions, this is sufficient to determine the metric in all of spacetime.

\subsubsection{Kruskal diagram}

To see where the event horizon is mapped in Kruskal spacetime:
\begin{equation*}
  r = 2GM
  \qquad \implies \qquad
  U = 0 \quad \lor \quad V = 0
\end{equation*}
This means that the event horizon no longer is a null surface, but two null surfaces intersecting at $ U = V = 0 $, in agreement with the near horizon limit (Rindler space): the null surface $ U = 0 $ is the horizon of the black hole, known as the \textit{future horizon}; the null surface $ V = 0 $ is the horizon of the white hole, known as the \textit{past horizon}.
On the other hand, the singularity becomes:
\begin{equation*}
  r = 0
  \qquad \implies \qquad
  UV = 1
\end{equation*}
This hyperbola has two disconnected components: one ($ U,V > 0 $) corresponding to the singularity of the black hole and one ($ U,V < 0 $) to the singularity of the white hole.

\begin{figure}
  \centering
  \includegraphics[width = 0.50 \textwidth]{kruskal.png}
  \caption{Kruskal diagram for the Schwarzschild balck hole: green lines are $ r = \text{const.} $ and blue lines $ t = \text{const.} $, while singularities are in red and horizons in black.}
  \label{fig:kruskal}
\end{figure}

All of this information can be represented in a \bctxt{Kruskal diagram}, as shown in \figref{fig:kruskal}: the $ U $ and $ V $ axes are at 45°, reflecting that they are null lines, and these are the two horizons. In this diagram, the vertical direction can be viewed as a time coordinate $ T = \frac{1}{2} (V + U) $ and the horizontal one as a spatial coordinate $ X = \frac{1}{2} (V - U) $. This diagram makes it clear that the black hole and the white hole cohabit the same spacetime. Lines of constant $ r $ and $ t $ correspond respectively to lines of constant $ UV $ and $ U/V $, as from \eeref{eq:kruskal-coord-1}{eq:kruskal-coord-2}.
The Kruskal diagram makes it clear that, once an observer crosses the event horizon, the singularity is unavoidable: the $ r = 0 $ are not timelike worldlines, but \textit{the singularity is spacelike}. Therefore, for the in-falling observer, after crossing the horizon the singularity of the black hole does not lie in a spatial direction, but it lies in the future. Similarly, the singularity of the white hole lies in the past. This can be formally seen using the Killing vector field $ K = \pa_t $ of the Schwarzschild solution. Outside of the horizon $ K $ is timelike, therefore it defines the conserved energy of geodesics; in Kruskal--Szekeres coordinates:
\begin{equation*}
  K = \pa_t = \frac{\pa V}{\pa t} \pa_V + \frac{\pa U}{\pa t} \pa_U = \frac{1}{4GM} \left( V \frac{\pa}{\pa V} - U \frac{\pa}{\pa U} \right)
  \quad \implies \quad
  g_{\mu \nu} K^\mu K^\nu = - \left( 1 - \frac{2GM}{r} \right)
\end{equation*}
Ouside of the horizon $ K_\mu K^\mu < 0 $ as expected. But inside, with $ r < 2GM $, the Killing vector field is spacelike: therefore, by \dref{def:stationary-spacetime} and \pref{prop:time-independent-metric}, the full Schwarzschild spacetime is not time-independent, but it becomes clear only after crossing the event horizon.

\paragraph{Einstein--Rosen bridge}

With reference to the Kruskal diagram in \figref{fig:kruskal}, the right-hand quadrant is the exterior of the black hole, covered by the Schwarzschild coordinates, while the upper and lower quadrants are respectively the interior of the black hole and of the white hole, covered by ingoing and outgoing Eddington-Finkelstein coordinates (see \figref{fig:kruskal-rind}). There remains the left-hand quadrant: this turns out to be just another copy of the exterior of the black hole, now covered by $ U > 0, V < 0 $.

\begin{figure}
  \centering
  \includegraphics[width = 0.70 \textwidth]{kruskal-rind.png}
  \caption{Regions of Kruskal spacetime covered by ingoing and outgoing Eddington--Finkelstein coordinates respectively.}
  \label{fig:kruskal-rind}
\end{figure}

The full spacetime hence contains two asymptotically-flat regions, joined together by a black hole: note that it is not possible for an observer in one region to send a signal in the other, because the causal structure of the spacetime does not allow this.

To elucidate the spatial geometry connecting these two regions, one need to study the $ t = 0 $ hypersurface of Kruskal spacetime: this is a straight horizontal line passing through $ U = V = 0 $ in \figref{fig:kruskal}. In Schwarzschild coordinates, the geometry of the $ t = 0 $ hypersurface is given by:
\begin{equation}
  \dd s^2 = \left( 1 - \frac{2GM}{r} \right)^{-1} \dd r^2 + r^2 \left( \dd \vartheta^2 + \sin^2 \vartheta \, \dd \varphi^2 \right)
  \label{eq:wormhole-metric}
\end{equation}
which is valid for $ r > 2GM $. Kruskal spacetime presents two copies of this geometry, one in the right-hand quadrant and one in the left-hand one, which are glues together at $ r = 2GM $ to give a wormhole-like geometry as in \figref{fig:wormhole}: this is called the \bctxt{Einstein--Rosen bridge}. Note that it is not possible to travel through it, as the paths are spacelike, not timelike.

\begin{figure}
  \centering
  \includegraphics[width = 0.40 \textwidth]{wormhole.png}
  \includegraphics[width = 0.50 \textwidth]{wormhole-rho.png}
  \caption{The Einstein--Rosen bridge and the $ \rho $ coordiante which parametrizes it.}
  \label{fig:wormhole}
\end{figure}

It is possible to write a metric which includes both sides. Define the radial coordinate $ \rho $ such that:
\begin{equation*}
  r = \rho \left( 1 + \frac{GM}{2\rho} \right)^2 = \rho + GM + \frac{G^2 M^2}{4 \rho}
\end{equation*}
Plotting $ r = r(\rho) $, as in \figref{fig:wormhole}, it is clear that there are two values of $ \rho $ for each value of $ r > 2GM $, while $ r = 2GM $ only corresponds to $ \rho = GM/2 $: thus, is possible to parametrize one side of the wormhole with $ \rho \in (0, GM/2) $ and the other with $ \rho \in (GM/2, \infty) $. The metric in \eref{eq:wormhole-metric} becomes:
\begin{equation}
  \dd s^2 = \left( 1 + \frac{GM}{2\rho} \right)^4 \left[ \dd \rho^2 + \rho^2 \left( \dd \vartheta^2 + \sin^2 \vartheta \, \dd \varphi^2 \right) \right]
\end{equation}
It is clear that, as $ \rho \rightarrow \infty $, the (spatial) metric is that of $ \R^3 $. But the same is true for $ \rho \rightarrow 0 $: note that there's a symmetry $ r \mapsto r $ for $ \rho \mapsto G^2 M^2 / (4\rho) $, which swaps the two asymptotic spacetimes while leaving the midpoint at $ \rho = GM/2 $ invariant, i.e. it swaps the two sides of the wormhole.

Furthermore, the radius of the spatial $ \Sn{2} $ is minimum at $ 2GM $ in the midpoint $ \rho = GM/2 $ of the wormhole and increases in either direction: this middle point corresponds to the common point $ U = V = 0 $ of the two horizons and is called the \bctxt{bifurcation sphere}.

Although there is no way that an observer in one quadrant can signal an observer in the other quadrant, there is one way for them to communicate: jump into the black hole and meet behind the respective horizons. However, as the white hole is thought to have no physical manifestation, the other universe in the left-hand quadrant of Kruskal space is thought to be a mathematical artifact too.

\subsubsection{Penrose diagram}
\label{sssec:bh-schwarzshild-conformal-diag}

\begin{figure}
  \centering
  \includegraphics[width = 0.70 \textwidth]{kruskal-penrose.png}
  \caption{The Penrose diagram for Kruskal spacetime.}
  \label{fig:kru-penr}
\end{figure}

From the Kruskal diagram in \figref{fig:kruskal}, it is easy to draw the Penrose diagram for Kruskal spacetime.
First, introduce new coordinates which cover the entire spacetime in a finite range:
\begin{equation*}
  U = \tan \tilde{U}
  \qquad \qquad
  V = \tan \tilde{V}
\end{equation*}
The new coordinates have range $ (\tilde{U},\tilde{V}) \in (-\frac{\pi}{2}, + \frac{\pi}{2}) $ and the Kruskal metric becomes:
\begin{equation*}
  \dd s^2 = \frac{1}{\cos^2 \tilde{U} \cos^2 \tilde{V}} \left[ - \frac{32 (GM)^3}{r} e^{-\frac{r}{2GM}} \dd \tilde{U}\,\dd \tilde{V} + r^2 \cos^2 \tilde{U} \cos^2 \tilde{V} \, \dd \Omega_2^2 \right]
\end{equation*}
This metric is conformally equivalent to:
\begin{equation*}
  \dd s^2 = - \frac{32 (GM)^3}{r} e^{-\frac{r}{2GM}} \dd \tilde{U}\,\dd \tilde{V} + r^2 \cos^2 \tilde{U} \cos^2 \tilde{V} \, \dd \Omega_2^2
\end{equation*}
However, the singularity at $ r = 0 $ must be retained: this is mapped to $ UV = 1 $, thus in the finite-range coordinates:
\begin{multline*}
  \tan \tilde{U} \tan \tilde{V} = 1
  \quad \iff \quad
  \sin \tilde{U} \sin \tilde{V} - \cos \tilde{U} \cos \tilde{V} = 0 \\
  \quad \iff \quad
  \cos (\tilde{U} + \tilde{V}) = 0
  \quad \iff \quad
  \tilde{U} + \tilde{V} = \pm \frac{\pi}{2}
\end{multline*}
These are straight horizontal lines in the Penrose diagram which chop the diamond-shaped diagram: the resulting diagram is shown in \figref{fig:kru-penr}, with singularities represented by jagged lines.

The Penrose diagram contains the same information as the Kruskal diagram in \figref{fig:kruskal}, but it allows to state a more rigorous definition of black hole. Restricting to aymptotically flat spacetimes, the two asymptotic regions contain both null infinities $ \mathcal{I}^+ $ and $ \mathcal{I}^- $: the black hole region is then defined as the set of points which cannot send signals to $ \mathcal{I}^+ $. The boundary of the black hole region is the \textit{future event horizon} $ \mathcal{H}^+ $: equivalently, $ \mathcal{H}^+ $ is the boundary of the causal past of $ \mathcal{I}^+ $. With reference to \figref{fig:kru-penr}, the black hole associated to $ \mathcal{I}^+ $ is the upper quadrant and the left quadrant, while the black hole associated to $ \mathcal{I}'^{+} $ is the upper quadrant and the right quadrant.

Importantly, note that to define a black hole one needs to know the whole spacetime, as all lightrays must be run backwards from $ \mathcal{I}^+ $ to determine $ \mathcal{H}^+ $ as their boundary: there is no reference to any spacelike hypersurface at constant coordinate time, thus an observer cannot really know if it is inside a black hole unless by knowing the entire future evolution of the spacetime.

Equivalently, it is possible to define the white hole region and the \textit{past event horizon} as the region that cannot receive signals from $ \mathcal{I}^- $ and its boundary.

\subsection{Weak cosmic censorship}

Kruskal spacetime is unphysical in a number of ways. In fact, in reality, black holes do not emerge from white holes, but are formed by collapsing stars: the resulting Penrose diagram is rather different from that in \figref{fig:kru-penr}. To understand the causal structure of such a black hole, consider the (unrealistic) situation of the spherically-symmetric collapse of a shell of matter: inside the shell, spacetime is flat, meanwhile on the outside it is described by Schwarzschild metric \eref{eq:schw-metric} (by Birkhoff's theorem, time dependence does not change this fact). Furthermore, assume (unrealistically) that the shell is travelling at the speed of light, so that the Penrose diagrams for Minkowski spacetime and Kruskal spacetime can be glued as in \figref{fig:cens-1}: this is the Penrose diagram of a collapsing black hole.

\begin{figure}
  \centering
  \includegraphics[width = 0.70 \textwidth]{cens-1.png}
  \caption{Joining two Penrose diagrams.}
  \label{fig:cens-1}
\end{figure}

\begin{figure}[b]
  \centering
  \includegraphics[width = 0.20 \textwidth]{cens-2.png}
  \caption{Penrose diagram for a collapsing star forming a black hole.}
  \label{fig:cens-2}
\end{figure}

Despite the number of unphysical assumptions, the Penrose diagram thus derived correctly describes the causal structure of a realistic spherically-symmetric collapsing star: as in \eref{eq:kruskal-metric}, its surface follows a timelike trajectory.

\subsubsection{Naked singularities}

An important feature of the black hole is that the singularity remains shrouded behind the event horizon, so that an asymptotic observer cannot see it.

On the contrary, Einstein field equations also allow \bctxt{naked singularities}, i.e. singularities without an horizon. For example, consider Schwarzschild metric with $ M < 0 $:
\begin{equation*}
  ds^2 = - \left( 1 + \frac{2G \abs{M}}{r} \right) dt^2 + \left( 1 + \frac{2G \abs{M}}{r} \right)^{-1} dr^2 + r^2 \left( d\vartheta^2 + \sin^2 \vartheta \, d\varphi^2 \right)
\end{equation*}
This metric presents no coordinate singularity at $ r = 2G \abs{M} $, and so no event horizon. The Penrose diagram for such a spacetime is found analogously to that of Minkwoski spacetime, setting null coordinates $ u = t - r $ and $ v = t + r $. The resulting diagram, shown in \figref{fig:cens-3}, is identical to that of Minkowski spacetime (\figref{fig:conformal-minkowski}), with the difference that now $ r = 0 $ presents a curvature singularity not shielded by an horizon: it is thus observed from $ \mathcal{I}^+ $.

\begin{figure}
  \centering
  \includegraphics[width = 0.70 \textwidth]{cens-3.png}
  \caption{Penrose diagram for a negative mass black hole and unphysical collapse of a star.}
  \label{fig:cens-3}
\end{figure}

An important conjecture in General Relativity, known as the \bctxt{weak cosmic censorship conjecture}, states the following: given matter which obeys the dominant energy condition \eref{eq:dominant-en-cond}, generic smooth initial conditions on a spatial hypersurface for both the metric and matter fields in an asymptotically-flat spacetime will not evolve to form naked singularities.

If this conjecture is true, then a dynamical evolution such as that in \figref{fig:cens-3} must be ruled out: in this diagram, once the singularity forms, fields can no longer be evolved beyond the shown lightray; strictly speaking, this means that the dynamical evolution stops at the red line and cannot be extended further, thus abruptly ending $ \mathcal{I}^+ $.

There is no proof of the weak cosmic censorship conjecture, but only circumstantial evidence. However, there is one naked singularity that seems to be physical: the Big Bang singularity. In fact, since it is in the far past, it does not violate cosmic censorship.

\subsection{de Sitter black holes}

Schwarzschild metric \eref{eq:schw-metric} solves the Einstein field equations with $ \Lambda = 0 $, i.e. in asymptotically Minkowski spacetime. To generalize, consider $ \Lambda \neq 0 $, so that the field equations become $ R_{\mu \nu} = \Lambda g_{\mu \nu} $, and use the following (already seen) ansatz:
\begin{equation*}
  \dd s^2 = - f(r)^2 \dd t^2 + f(r)^{-2} \dd r^2 + r^2 \left( \dd \vartheta^2 + \sin^2 \vartheta \, \dd \varphi^2 \right)
\end{equation*}
The field equations imply that:
\begin{equation*}
  f'' + \frac{2f'}{r} + \frac{{f'}^2}{f} = - \frac{\Lambda}{f}
  \qquad \land \qquad
  1 - 2 f f' r - f^2 = \Lambda r^2
\end{equation*}
The most general solution is:
\begin{equation}
  f(r)^2 = 1 - \frac{2GM}{r} \mp \frac{r^2}{R^2}
  \qquad
  R^2 = \frac{3}{\abs{\Lambda}}
  \label{eq:6.14}
\end{equation}
where the negative sign is the solution with $ \Lambda > 0 $ and the positive sign with $ \Lambda < 0 $, hence describing black holes in de Sitter and anti-de Sidder spacetime respectively: in fact, if $ 2GM \ll R^2 $, then for $ r \ll R $ the metric is that of a Schwarzschild black hole in Minkowski spacetime.

\section{Reissner--Nordström black holes}
\label{sec:bh-rn}

\section{Kerr black holes}
\label{sec:bh-kerr}

\section{Dynamical black holes}
\label{sec:bh-dynamics}










